\documentclass[10pt,letterpaper]{article}
\usepackage[letterpaper,margin=0.72in,headheight=16pt,footskip=28pt]{geometry}
\usepackage[T1]{fontenc}
\usepackage{lmodern}
\usepackage{microtype}
\usepackage[dvipsnames,table]{xcolor}
\usepackage{graphicx}
\usepackage{booktabs}
\usepackage{tabularx}
\usepackage{longtable}
\usepackage{array}
\usepackage{ragged2e}
\usepackage{pdflscape}
\usepackage{enumitem}
\usepackage{parskip}
\usepackage{needspace}
\usepackage[most]{tcolorbox}
\usepackage{listings}
\usepackage{upquote}
\usepackage{textcomp}
\usepackage{fancyhdr}
\usepackage{titlesec}
\usepackage{lastpage}
\usepackage{xurl}
\usepackage[hidelinks]{hyperref}
\usepackage{bookmark}

\definecolor{Navy}{HTML}{123047}
\definecolor{Teal}{HTML}{0E6670}
\definecolor{LightTeal}{HTML}{EAF4F4}
\definecolor{Critical}{HTML}{A12828}
\definecolor{High}{HTML}{A45D00}
\definecolor{Standard}{HTML}{356A45}
\definecolor{Defense}{HTML}{5D6470}
\definecolor{Ink}{HTML}{20262D}
\definecolor{Muted}{HTML}{66717B}
\definecolor{Panel}{HTML}{F5F7F8}
\definecolor{RuleLine}{HTML}{D7DEE3}

\hypersetup{
  pdftitle={GitHub Advanced Security Supported-Language Secure Coding Standards Catalog},
  pdfauthor={Application Security Engineering},
  colorlinks=true,
  linkcolor=Navy,
  urlcolor=Teal,
  citecolor=Navy
}

\pagestyle{fancy}
\fancyhf{}
\lhead{\textcolor{Muted}{GHAS Secure Coding Standards}}
\rhead{\textcolor{Muted}{Enterprise Baseline v1.0}}
\cfoot{\textcolor{Muted}{\thepage\ / \pageref*{LastPage}}}
\renewcommand{\headrulewidth}{0.3pt}
\renewcommand{\footrulewidth}{0pt}

\titleformat{\section}{\Large\bfseries\color{Navy}}{\thesection}{0.65em}{}
\titleformat{\subsection}{\large\bfseries\color{Teal}}{\thesubsection}{0.6em}{}
\titleformat{\subsubsection}{\normalsize\bfseries\color{Ink}}{\thesubsubsection}{0.5em}{}
\titlespacing*{\section}{0pt}{1.7em}{0.7em}
\titlespacing*{\subsection}{0pt}{1.3em}{0.5em}
\setcounter{tocdepth}{2}
\setcounter{secnumdepth}{2}

\newcommand{\code}[1]{\texttt{\small #1}}
\newcommand{\prioritybadge}[2]{\tcbox[on line,boxrule=0pt,arc=1.2mm,boxsep=0.7mm,left=1mm,right=1mm,top=0.4mm,bottom=0.4mm,colback=#2!12,coltext=#2]{\scriptsize\bfseries #1}}
\newcommand{\pill}[1]{\tcbox[on line,boxrule=0pt,arc=1mm,boxsep=0.6mm,left=0.8mm,right=0.8mm,top=0.3mm,bottom=0.3mm,colback=LightTeal,coltext=Teal]{\scriptsize\bfseries #1}}

\newtcolorbox{rulebox}[1][]{
  enhanced, breakable, colback=white, colframe=RuleLine,
  boxrule=0.55pt, arc=1.4mm, left=2.5mm,right=2.5mm,top=1.5mm,bottom=1.5mm,
  borderline west={2.2pt}{0pt}{Teal},
  fonttitle=\bfseries\color{Navy}, coltitle=Navy,
  colbacktitle=Panel, attach boxed title to top left={xshift=1.5mm,yshift*=-1mm},
  boxed title style={boxrule=0pt,arc=1mm,colback=Panel},
  before skip=8pt, after skip=8pt, #1
}

\newtcolorbox{callout}[2][]{
  enhanced, breakable, colback=#2!5, colframe=#2!40,
  boxrule=0.6pt, arc=1.3mm, left=3mm,right=3mm,top=2mm,bottom=2mm,
  #1
}

\lstdefinestyle{securecode}{
  basicstyle=\ttfamily\footnotesize,
  columns=fullflexible,
  keepspaces=true,
  showstringspaces=false,
  breaklines=true,
  breakatwhitespace=false,
  frame=single,
  rulecolor=\color{RuleLine},
  backgroundcolor=\color{Panel},
  xleftmargin=2mm,
  xrightmargin=1mm,
  aboveskip=6pt,
  belowskip=8pt,
  tabsize=2
}
\lstset{style=securecode}

\newcolumntype{Y}{>{\RaggedRight\arraybackslash}X}
\newcolumntype{P}[1]{>{\RaggedRight\arraybackslash}p{#1}}
\renewcommand{\arraystretch}{1.15}
\setlength{\emergencystretch}{3em}

\begin{document}
\hypersetup{pageanchor=false}
\begin{titlepage}
\thispagestyle{empty}
\vspace*{1.0in}
{\color{Teal}\rule{1.2in}{3pt}}\\[0.35in]
{\Huge\bfseries\color{Navy}GitHub Advanced Security\\[0.12in]Supported-Language\\[0.12in]Secure Coding Standards Catalog\par}
\vspace{0.28in}
{\Large\color{Teal}Enterprise Secure Coding \& CodeQL Enforcement Handout\par}
\vspace{0.55in}
\begin{callout}{Teal}
\textbf{Purpose.} Establish a traceable control system connecting language-level secure coding requirements to weakness classes, implementation patterns, CodeQL detection, and SSDLC enforcement.
\end{callout}
\vfill
\begin{tabular}{@{}ll@{}}
\textbf{Baseline} & v1.0 \\
\textbf{Verified} & 23 September 2026 \\
\textbf{Scope} & CodeQL-supported languages + GitHub Actions workflow analysis \\
\textbf{Source policy} & Official and authoritative primary sources only \\
\textbf{Normative terms} & MUST, MUST NOT, SHOULD, SHOULD NOT, MAY \\
\end{tabular}
\vspace{0.7in}
\end{titlepage}
\hypersetup{pageanchor=true}
\pagenumbering{roman}
\tableofcontents
\clearpage
\pagenumbering{arabic}

\section{Executive Overview}
This catalog establishes a traceable control system from \textbf{secure-coding requirement} $\rightarrow$ \textbf{weakness} $\rightarrow$ \textbf{implementation behavior} $\rightarrow$ \textbf{CodeQL or companion control} $\rightarrow$ \textbf{SSDLC enforcement}. It applies a shared 20-domain baseline to every supported language and then adds language-specific overlays where the language or runtime has distinct security semantics.

\begin{callout}{Navy}
\textbf{Interpretation boundary.} CodeQL support is not evidence that a language has a complete secure-coding standard. Likewise, a clean CodeQL result is not evidence that an application is secure. Authentication, authorization, business logic, runtime configuration, secrets, dependency provenance, deployment architecture, and many cryptographic design decisions require controls beyond static analysis.
\end{callout}

\subsection{Research and Control Method}
\begin{enumerate}[leftmargin=*,itemsep=2pt]
  \item Verify the current CodeQL language, version, framework, and build-mode matrix from GitHub.
  \item Identify the strongest official secure-coding authority for each language.
  \item Establish one cross-language requirement for each requested security domain.
  \item Add language-specific overlays instead of copying generic rules into every chapter.
  \item Map rules to CWE only where the relationship is defensible.
  \item Record CodeQL coverage as direct, partial, not reliably detectable, or requiring another control.
  \item Use secret scanning, push protection, Dependabot, dependency review, and Actions/OIDC where CodeQL is not the primary enforcement mechanism.
  \item Preserve exact CodeQL query IDs only when verified against GitHub's official query help.
\end{enumerate}

\subsection{Priority Model}
\begin{tabularx}{\textwidth}{@{}P{1.45in}Y@{}}
\toprule
\textbf{Priority} & \textbf{Objective basis}\\
\midrule
\prioritybadge{Critical Control}{Critical} & Failure can directly enable remote code execution, memory corruption, authentication or authorization bypass, secret compromise, major injection, cryptographic trust failure, or supply-chain compromise.\\
\prioritybadge{High Priority}{High} & Failure can materially compromise confidentiality, integrity, or availability, or undermine a security boundary, but usually needs additional conditions or has narrower impact.\\
\prioritybadge{Standard Practice}{Standard} & Security-relevant engineering hygiene that reduces defect likelihood and improves verification.\\
\prioritybadge{Defense in Depth}{Defense} & Additional protection whose absence does not by itself establish a vulnerability but reduces resilience or blast radius.\\
\bottomrule
\end{tabularx}
\clearpage
\begin{landscape}
\section{Verified GHAS / CodeQL Language Matrix}
\textbf{Language matrix verified on: 23 September 2026.} Framework and library coverage summarizes explicitly documented modeling and is not a promise that every API is modeled.

\scriptsize
\setlength{\tabcolsep}{3pt}
\begin{longtable}{@{}P{0.72in}P{0.9in}P{1.45in}P{3.55in}P{1.25in}P{0.65in}@{}}
\toprule
\textbf{Language / target} & \textbf{CodeQL ID} & \textbf{Current versions} & \textbf{Framework / library coverage summary} & \textbf{Build mode} & \textbf{Compiled?}\\
\midrule
\endfirsthead
\toprule
\textbf{Language / target} & \textbf{CodeQL ID} & \textbf{Current versions} & \textbf{Framework / library coverage summary} & \textbf{Build mode} & \textbf{Compiled?}\\
\midrule
\endhead
C & c-cpp & C89, C99, C11, C17, C23 (C23 beta) & Berkeley sockets, string.h; extensible library modeling & none / autobuild / manual & Yes\\
C++ & c-cpp & C++98, 03, 11, 14, 17, 20, 23 (C++23 beta; C++20 modules unsupported) & C/C++ built-ins; extensible library modeling & none / autobuild / manual & Yes\\
C\# & csharp & C\# through 14; .NET Framework 4.8; .NET Core 3.1; .NET 5--10 & ASP.NET/Core/Razor, Dapper, EF/EF Core, Json.NET, NHibernate & none / autobuild / manual & Yes\\
Go & go & Go through 1.27 & net/http ecosystem; Gin, Echo, GORM, Gorilla, protobuf, common DB/log/network libs & autobuild / manual & Yes\\
Java & java-kotlin & Java 7--26; language features understood for Java 8--26 & JDBC/JPA, Spring MVC/JDBC, Jackson, JAX-RS, Hibernate, MyBatis, Struts, serialization libs & none / autobuild / manual & Yes\\
Kotlin & java-kotlin & Kotlin 1.8.0--2.4.20 & Java/Kotlin pack; JVM libraries/frameworks modeled by java-all & autobuild / manual; build required for Kotlin & Yes\\
JavaScript & javascript-typescript & ECMAScript 2022 or lower & Browser, Node, Express, Fastify, React, Vue, Angular, common DB/template/network libs & none & No\\
TypeScript & javascript-typescript & TypeScript 2.6--7.0 & Analyzed by JavaScript extractor with TypeScript enabled; same modeled ecosystem & none & No\\
Python & python & Python 2.7 and 3.5--3.14 & Django, Flask, FastAPI, aiohttp, requests/httpx, SQLAlchemy, PyYAML, cryptography, many DB libs & none & No\\
Ruby & ruby & Ruby through 3.3 & Ruby on Rails, HTTP clients, Nokogiri/libxml, rubyzip & none & No\\
Rust & rust & Rust editions 2021 and 2024; nightly features unsupported & std/core/alloc, Tokio, Actix/Rocket/Warp, Reqwest, Serde, SQLx, Rustls, common DB/crypto libs & none only (no-build extraction) & No build tracing\\
Swift & swift & Swift 5.4--6.3; Embedded Swift unsupported & Foundation, CryptoKit, Network, Core Data, WebKit, common DB/network/XML libs & autobuild / manual; macOS required & Yes\\
GitHub Actions & actions & Workflow and action-metadata YAML & \code{.github/workflows/*.yml|yaml}, \code{action.yml|yaml} & not applicable & No\\
\bottomrule
\end{longtable}
\normalsize
\begin{callout}{Teal}
\textbf{Current caveats.} The catalog records point-in-time support details and must be reverified when the CodeQL bundle changes. In particular, beta or unsupported language features and platform requirements can materially affect extraction and analysis.
\end{callout}
\end{landscape}
\begin{landscape}
\section{Secure Coding Standards Authority Matrix}
The classifications below deliberately distinguish a formal secure-coding standard from vendor guidance, official security guidance, and an enterprise baseline derived from multiple first-party sources.

\scriptsize
\setlength{\tabcolsep}{3pt}
\begin{longtable}{@{}P{0.8in}P{2.05in}P{1.45in}P{3.25in}P{0.7in}@{}}
\toprule
\textbf{Language / target} & \textbf{Primary authority} & \textbf{Classification} & \textbf{Secondary authorities} & \textbf{Source}\\
\midrule
\endfirsthead
\toprule
\textbf{Language / target} & \textbf{Primary authority} & \textbf{Classification} & \textbf{Secondary authorities} & \textbf{Source}\\
\midrule
\endhead
C & SEI CERT C Coding Standard (2016) & Formal Secure Coding Standard & SEI CERT + NIST SSDF + MITRE CWE + OWASP where web-facing & CERT-C\\
C++ & SEI CERT C++ Coding Standard (2016) & Formal Secure Coding Standard & SEI CERT + NIST SSDF + MITRE CWE + OWASP where web-facing & CERT-CPP\\
C\# & Microsoft .NET secure coding/security documentation & Vendor Secure Coding Guidelines & ASP.NET Core security docs + NIST/OWASP/MITRE + GitHub & DOTNET-SEC\\
Go & Go Security Best Practices and standard-library docs & Official Security Guidance & NIST/OWASP/MITRE + GitHub & GO-SEC\\
Java & SEI CERT Oracle Coding Standard for Java + Oracle Secure Coding Guidelines for Java SE & Formal Secure Coding Standard & Oracle current Java SE guidance + NIST/OWASP/MITRE + GitHub & CERT-JAVA\\
Kotlin & Kotlin security/null-safety docs + JVM guidance & Derived Secure Coding Baseline & Oracle Java guidance + NIST/OWASP/MITRE + GitHub & KOTLIN-SEC\\
JavaScript & Node.js security/API docs + ECMAScript/runtime semantics & Derived Secure Coding Baseline & OWASP ASVS + NIST/MITRE + GitHub; browser framework docs as applicable & NODE-SEC\\
TypeScript & TypeScript Handbook + TSConfig + JavaScript runtime baseline & Derived Secure Coding Baseline & Node/browser guidance + OWASP/NIST/MITRE + GitHub & TS-TYPES\\
Python & Python language/stdlib security warnings and API docs & Derived Secure Coding Baseline & NIST/OWASP/MITRE + GitHub & PY-PICKLE\\
Ruby & Ruby core/stdlib security docs + Rails Security Guide & Derived Secure Coding Baseline & NIST/OWASP/MITRE + GitHub & RUBY-MARSHAL\\
Rust & Rustonomicon + standard/Cargo documentation & Derived Secure Coding Baseline & NIST/OWASP/MITRE + GitHub & RUST-UNSAFE\\
Swift & Swift language safety docs + Apple platform security docs & Derived Secure Coding Baseline & NIST/OWASP/MITRE + GitHub & SWIFT-MEM\\
GitHub Actions & GitHub Secure Use Reference + CodeQL Actions queries & Official Security Guidance & NIST SSDF + GitHub supply-chain/secret/OIDC docs & GH-SECURE-ACTIONS\\
\bottomrule
\end{longtable}
\end{landscape}
\normalsize

\begin{itemize}[leftmargin=*,itemsep=2pt]
  \item \textbf{Formal Secure Coding Standard}: a named secure-coding standard is explicitly published by an authoritative standards or security body.
  \item \textbf{Vendor Secure Coding Guidelines}: the platform owner publishes a cohesive secure-coding guideline, but it is not represented here as a formal language standard.
  \item \textbf{Official Security Guidance}: the language/project publishes first-party security best practices or secure-use guidance.
  \item \textbf{Derived Secure Coding Baseline}: security rules are distributed across official language, runtime, API, and platform documents and are normalized here into one enterprise control set.
\end{itemize}
\clearpage
\section{Cross-Language Secure Coding Baseline}
These controls apply to every supported language unless a language chapter marks the domain as not applicable or narrows the requirement. Language-specific overlays strengthen the baseline; they do not weaken it.

\begin{rulebox}[title={\texttt{SCS-ALL-INPUT-001} \hfill \prioritybadge{High Priority}{High}}]
\textbf{Domain.} Input Validation

\textbf{Requirement.} All data crossing a trust boundary MUST be validated for expected type, syntax, length, range, structure, and business constraints before security-sensitive use.

\textbf{Rationale.} Type systems and parser success do not establish that attacker-controlled data is authorized or semantically safe.

\textbf{Risk.} Injection, parser abuse, resource exhaustion, authorization bypass

\begin{description}[leftmargin=2.65cm,labelwidth=2.3cm,style=nextline,itemsep=1pt,topsep=3pt]
  \item[CWE] CWE-20
  \item[OWASP] OWASP ASVS 5.0 --- validation/business logic
  \item[CERT] Not mapped
  \item[CodeQL] \textbf{Partial}; Language-specific taint/data-flow queries
  \item[Enforcement] CodeQL + schema/parser validation + unit/fuzz tests
  \item[Automation] Partial
  \item[Sources] \texttt{NIST-SSDF, OWASP-ASVS, MITRE-CWE, JAVA-SCG}
\end{description}
\end{rulebox}

\begin{rulebox}[title={\texttt{SCS-ALL-OUTPUT-001} \hfill \prioritybadge{High Priority}{High}}]
\textbf{Domain.} Output Encoding

\textbf{Requirement.} Untrusted data MUST be encoded or escaped for the exact output context; developers MUST prefer framework APIs that perform contextual escaping by default.

\textbf{Rationale.} Encoding is context-sensitive; HTML, JavaScript, URLs, headers, XML, SQL, and shell contexts are not interchangeable.

\textbf{Risk.} XSS, response splitting, code/command interpretation

\begin{description}[leftmargin=2.65cm,labelwidth=2.3cm,style=nextline,itemsep=1pt,topsep=3pt]
  \item[CWE] CWE-79, CWE-116
  \item[OWASP] OWASP ASVS 5.0 --- encoding/sanitization
  \item[CERT] Not mapped
  \item[CodeQL] \textbf{Partial}; Language/framework-specific XSS and response-splitting queries
  \item[Enforcement] CodeQL + framework-safe templates + tests
  \item[Automation] Partial
  \item[Sources] \texttt{OWASP-ASVS, MITRE-CWE, GH-BUILTIN}
\end{description}
\end{rulebox}

\begin{rulebox}[title={\texttt{SCS-ALL-INJ-001} \hfill \prioritybadge{Critical Control}{Critical}}]
\textbf{Domain.} Injection Prevention

\textbf{Requirement.} Untrusted data MUST NOT be concatenated into SQL, OS commands, code, templates, expressions, or other interpreters; parameterized or structured APIs MUST be used where available.

\textbf{Rationale.} Separating code from data removes attacker control over interpreter syntax.

\textbf{Risk.} SQL/NoSQL injection, command injection, code injection, template/expression injection

\begin{description}[leftmargin=2.65cm,labelwidth=2.3cm,style=nextline,itemsep=1pt,topsep=3pt]
  \item[CWE] CWE-78, CWE-89, CWE-94, CWE-95, CWE-943
  \item[OWASP] OWASP ASVS 5.0 --- injection prevention
  \item[CERT] Not mapped
  \item[CodeQL] \textbf{Direct}; Representative: cpp/sql-injection; cs/sql-injection; go/sql-injection; java/sql-injection; js/sql-injection; py/sql-injection; rb/sql-injection; rust/sql-injection; swift/sql-injection
  \item[Enforcement] CodeQL required check + API-level parameterization
  \item[Automation] Yes
  \item[Sources] \texttt{GH-BUILTIN, MITRE-CWE, OWASP-ASVS}
\end{description}
\end{rulebox}

\begin{rulebox}[title={\texttt{SCS-ALL-MEM-001} \hfill \prioritybadge{Critical Control}{Critical}}]
\textbf{Domain.} Memory Safety

\textbf{Requirement.} Code that can perform unchecked memory access MUST establish bounds, initialization, lifetime, ownership, and aliasing invariants before access; memory-safe language code SHOULD remain inside safe subsets.

\textbf{Rationale.} Memory corruption can become arbitrary code execution or information disclosure.

\textbf{Risk.} Out-of-bounds access, use-after-free, double free, invalid pointer use

\begin{description}[leftmargin=2.65cm,labelwidth=2.3cm,style=nextline,itemsep=1pt,topsep=3pt]
  \item[CWE] CWE-119, CWE-416, CWE-415, CWE-825
  \item[OWASP] Context-dependent / not forced
  \item[CERT] SEI CERT C/C++ memory rules
  \item[CodeQL] \textbf{Partial}; Representative: cpp/use-after-free; rust/access-after-lifetime-ended
  \item[Enforcement] CodeQL + compiler/runtime checks + sanitizers/dynamic tests + review of unsafe/native boundaries
  \item[Automation] Partial
  \item[Sources] \texttt{CERT-C, CERT-CPP, RUST-UNSAFE, SWIFT-MEM, GH-BUILTIN}
\end{description}
\end{rulebox}

\begin{rulebox}[title={\texttt{SCS-ALL-AUTHN-001} \hfill \prioritybadge{Critical Control}{Critical}}]
\textbf{Domain.} Authentication

\textbf{Requirement.} Applications MUST use a vetted platform/framework/identity-provider authentication mechanism and MUST NOT implement custom credential cryptography or authentication protocols without specialized review.

\textbf{Rationale.} Authentication implementations are protocol- and state-sensitive and commonly fail at credential, session, recovery, or trust boundaries.

\textbf{Risk.} Authentication bypass, credential compromise

\begin{description}[leftmargin=2.65cm,labelwidth=2.3cm,style=nextline,itemsep=1pt,topsep=3pt]
  \item[CWE] CWE-287, CWE-306
  \item[OWASP] OWASP ASVS 5.0 --- authentication
  \item[CERT] Not mapped
  \item[CodeQL] \textbf{Partial}; Framework-specific authentication queries
  \item[Enforcement] Architecture review + framework configuration tests + DAST/IAST where applicable
  \item[Automation] Partial
  \item[Sources] \texttt{OWASP-ASVS, NIST-SSDF, MITRE-CWE}
\end{description}
\end{rulebox}

\begin{rulebox}[title={\texttt{SCS-ALL-AUTHZ-001} \hfill \prioritybadge{Critical Control}{Critical}}]
\textbf{Domain.} Authorization

\textbf{Requirement.} Authorization MUST be enforced server-side at each protected operation and object boundary, SHOULD default deny, and MUST be based on trusted identity/authorization state rather than client-supplied claims alone.

\textbf{Rationale.} Authenticated users can still request objects or actions they are not permitted to access.

\textbf{Risk.} Missing authorization, IDOR/BOLA, privilege escalation

\begin{description}[leftmargin=2.65cm,labelwidth=2.3cm,style=nextline,itemsep=1pt,topsep=3pt]
  \item[CWE] CWE-285, CWE-862, CWE-863, CWE-639
  \item[OWASP] OWASP ASVS 5.0 --- access control
  \item[CERT] Not mapped
  \item[CodeQL] \textbf{Partial}; Language/framework-specific missing authorization / IDOR queries
  \item[Enforcement] CodeQL where modeled + authorization tests + architecture/code review
  \item[Automation] Partial
  \item[Sources] \texttt{OWASP-ASVS, MITRE-CWE, GH-BUILTIN}
\end{description}
\end{rulebox}

\begin{rulebox}[title={\texttt{SCS-ALL-CRYPTO-001} \hfill \prioritybadge{Critical Control}{Critical}}]
\textbf{Domain.} Cryptography

\textbf{Requirement.} Applications MUST use current platform cryptographic APIs and cryptographically secure random generators for security-sensitive values; custom cryptographic primitives and obsolete/broken algorithms MUST NOT be used.

\textbf{Rationale.} Cryptographic strength depends on algorithm, mode, parameters, randomness, key management, and API usage.

\textbf{Risk.} Predictable tokens, weak encryption/signatures, key compromise

\begin{description}[leftmargin=2.65cm,labelwidth=2.3cm,style=nextline,itemsep=1pt,topsep=3pt]
  \item[CWE] CWE-327, CWE-326, CWE-338
  \item[OWASP] OWASP ASVS 5.0 --- cryptography
  \item[CERT] Not mapped
  \item[CodeQL] \textbf{Partial}; Weak-crypto / insecure-randomness / TLS queries vary by language
  \item[Enforcement] CodeQL + platform API policy + cryptographic design review
  \item[Automation] Partial
  \item[Sources] \texttt{OWASP-ASVS, MITRE-CWE, DOTNET-CRYPTO, GO-RAND, JAVA-26, APPLE-CRYPTO}
\end{description}
\end{rulebox}

\begin{rulebox}[title={\texttt{SCS-ALL-SECRET-001} \hfill \prioritybadge{Critical Control}{Critical}}]
\textbf{Domain.} Secrets

\textbf{Requirement.} Credentials, API tokens, private keys, and other secrets MUST NOT be hard-coded in source or written to logs; repositories SHOULD enable GitHub secret scanning and push protection for supported secret types.

\textbf{Rationale.} Source history and logs are durable distribution channels; exposed credentials can be replayed.

\textbf{Risk.} Credential exposure and unauthorized access

\begin{description}[leftmargin=2.65cm,labelwidth=2.3cm,style=nextline,itemsep=1pt,topsep=3pt]
  \item[CWE] CWE-798, CWE-532, CWE-312
  \item[OWASP] OWASP ASVS 5.0 --- secrets/data protection
  \item[CERT] Not mapped
  \item[CodeQL] \textbf{Requires another control}; GitHub secret scanning / push protection rather than CodeQL for primary enforcement
  \item[Enforcement] Secret scanning + push protection + approved secret store + rotation process
  \item[Automation] Yes
  \item[Sources] \texttt{GH-SECRET, GH-PUSH, MITRE-CWE, NIST-SSDF}
\end{description}
\end{rulebox}

\begin{rulebox}[title={\texttt{SCS-ALL-ERROR-001} \hfill \prioritybadge{High Priority}{High}}]
\textbf{Domain.} Error / Exception Handling

\textbf{Requirement.} Security failures MUST fail closed; externally visible errors MUST NOT disclose secrets, credentials, stack traces, filesystem paths, internal queries, or other sensitive implementation details.

\textbf{Rationale.} Error handling is part of the security boundary and can otherwise disclose exploitable state or bypass checks.

\textbf{Risk.} Information exposure, fail-open authorization/authentication behavior

\begin{description}[leftmargin=2.65cm,labelwidth=2.3cm,style=nextline,itemsep=1pt,topsep=3pt]
  \item[CWE] CWE-209, CWE-703
  \item[OWASP] OWASP ASVS 5.0 --- error handling/logging
  \item[CERT] Not mapped
  \item[CodeQL] \textbf{Partial}; Information-disclosure and unchecked-error queries vary
  \item[Enforcement] Code review + negative-path tests + logging policy + CodeQL where available
  \item[Automation] Partial
  \item[Sources] \texttt{OWASP-ASVS, MITRE-CWE, NIST-SSDF}
\end{description}
\end{rulebox}

\begin{rulebox}[title={\texttt{SCS-ALL-LOG-001} \hfill \prioritybadge{High Priority}{High}}]
\textbf{Domain.} Logging / Monitoring

\textbf{Requirement.} Security-relevant events SHOULD be logged with structured fields; untrusted data MUST be encoded for the logging sink and secrets/sensitive payloads MUST NOT be logged.

\textbf{Rationale.} Logs are both security evidence and attacker-controlled sinks; injection or sensitive data exposure weakens monitoring and confidentiality.

\textbf{Risk.} Log injection, sensitive data exposure, audit gaps

\begin{description}[leftmargin=2.65cm,labelwidth=2.3cm,style=nextline,itemsep=1pt,topsep=3pt]
  \item[CWE] CWE-117, CWE-532
  \item[OWASP] OWASP ASVS 5.0 --- logging
  \item[CERT] Not mapped
  \item[CodeQL] \textbf{Partial}; Log-injection / cleartext-logging queries exist for multiple languages
  \item[Enforcement] CodeQL + centralized logging policy + tests
  \item[Automation] Partial
  \item[Sources] \texttt{OWASP-ASVS, MITRE-CWE, GH-BUILTIN}
\end{description}
\end{rulebox}

\begin{rulebox}[title={\texttt{SCS-ALL-FILE-001} \hfill \prioritybadge{High Priority}{High}}]
\textbf{Domain.} File / Path Security

\textbf{Requirement.} File operations MUST constrain attacker-controlled paths to an intended root, avoid unsafe symbolic-link/TOCTOU assumptions, use secure temporary-file APIs, and validate archive extraction destinations.

\textbf{Rationale.} Path normalization alone is not an authorization decision and filesystem state can change between checks and use.

\textbf{Risk.} Path traversal, arbitrary file read/write, symlink race, unsafe temp files, Zip Slip

\begin{description}[leftmargin=2.65cm,labelwidth=2.3cm,style=nextline,itemsep=1pt,topsep=3pt]
  \item[CWE] CWE-22, CWE-59, CWE-367, CWE-377
  \item[OWASP] OWASP ASVS 5.0 --- file handling
  \item[CERT] Not mapped
  \item[CodeQL] \textbf{Direct/Partial}; Path-injection/traversal, archive extraction, temporary-file, and TOCTOU queries vary by language
  \item[Enforcement] CodeQL + filesystem permission tests + archive policy
  \item[Automation] Partial
  \item[Sources] \texttt{GH-BUILTIN, MITRE-CWE, PY-TAR, NIST-SSDF}
\end{description}
\end{rulebox}

\begin{rulebox}[title={\texttt{SCS-ALL-DESER-001} \hfill \prioritybadge{Critical Control}{Critical}}]
\textbf{Domain.} Serialization / Deserialization

\textbf{Requirement.} Untrusted input MUST NOT be deserialized by APIs that can instantiate arbitrary types or invoke attacker-influenced code; data-only formats and restrictive/safe loaders MUST be preferred.

\textbf{Rationale.} Powerful object deserializers can convert data parsing into code execution.

\textbf{Risk.} Remote code execution, object injection, denial of service

\begin{description}[leftmargin=2.65cm,labelwidth=2.3cm,style=nextline,itemsep=1pt,topsep=3pt]
  \item[CWE] CWE-502
  \item[OWASP] OWASP ASVS 5.0 --- deserialization
  \item[CERT] Not mapped
  \item[CodeQL] \textbf{Direct/Partial}; Representative: cs/unsafe-deserialization-untrusted-input; java/unsafe-deserialization; py/unsafe-deserialization; rb/unsafe-deserialization
  \item[Enforcement] CodeQL + API allowlist + schema validation
  \item[Automation] Yes/Partial
  \item[Sources] \texttt{DOTNET-BF, PY-PICKLE, RUBY-MARSHAL, JAVA-SCG, GH-BUILTIN}
\end{description}
\end{rulebox}

\begin{rulebox}[title={\texttt{SCS-ALL-NET-001} \hfill \prioritybadge{Critical Control}{Critical}}]
\textbf{Domain.} Network Communications

\textbf{Requirement.} Sensitive or authenticated network traffic MUST use TLS with normal certificate and hostname verification; production code MUST NOT disable or bypass trust validation except through a documented, narrowly scoped exception.

\textbf{Rationale.} Encryption without peer authentication remains vulnerable to interception and impersonation.

\textbf{Risk.} MITM, cleartext disclosure, downgrade

\begin{description}[leftmargin=2.65cm,labelwidth=2.3cm,style=nextline,itemsep=1pt,topsep=3pt]
  \item[CWE] CWE-295, CWE-297, CWE-319, CWE-757
  \item[OWASP] OWASP ASVS 5.0 --- communications
  \item[CERT] Not mapped
  \item[CodeQL] \textbf{Direct/Partial}; Representative: go/disabled-certificate-check; swift/insecure-tls; TLS/certificate queries in other packs
  \item[Enforcement] CodeQL + platform network policy + integration tests
  \item[Automation] Partial
  \item[Sources] \texttt{OWASP-ASVS, MITRE-CWE, APPLE-ATS, GH-BUILTIN}
\end{description}
\end{rulebox}

\begin{rulebox}[title={\texttt{SCS-ALL-WEB-001} \hfill \prioritybadge{Critical Control}{Critical}}]
\textbf{Domain.} Web Application Security

\textbf{Requirement.} Web applications MUST use framework protections for contextual output encoding, anti-CSRF where browser credentials are ambient, secure cookies/session handling, constrained CORS, validated redirects, and outbound-request controls against SSRF.

\textbf{Rationale.} Web security depends on browser/proxy semantics and trust boundaries not provided by the base programming language.

\textbf{Risk.} XSS, CSRF, SSRF, open redirect, session compromise

\begin{description}[leftmargin=2.65cm,labelwidth=2.3cm,style=nextline,itemsep=1pt,topsep=3pt]
  \item[CWE] CWE-79, CWE-352, CWE-918, CWE-601
  \item[OWASP] OWASP ASVS 5.0 --- web/browser/session controls
  \item[CERT] Not mapped
  \item[CodeQL] \textbf{Direct/Partial}; Representative: go/request-forgery; js/request-forgery; rust/request-forgery; framework-specific XSS/CSRF queries
  \item[Enforcement] CodeQL + framework config tests + DAST/IAST + architecture review
  \item[Automation] Partial
  \item[Sources] \texttt{OWASP-ASVS, GH-BUILTIN, ASPNET-CSRF, RAILS-SEC}
\end{description}
\end{rulebox}

\begin{rulebox}[title={\texttt{SCS-ALL-DEP-001} \hfill \prioritybadge{Critical Control}{Critical}}]
\textbf{Domain.} Dependency / Package Security

\textbf{Requirement.} Dependencies MUST come from approved sources, SHOULD be locked/pinned reproducibly, MUST be reviewed for known vulnerabilities, and MUST be updated or mitigated according to risk; dependency-review gates SHOULD protect pull requests where available.

\textbf{Rationale.} A secure application can be compromised by a vulnerable or malicious transitive dependency.

\textbf{Risk.} Known-vulnerable component, dependency confusion, malicious package/action

\begin{description}[leftmargin=2.65cm,labelwidth=2.3cm,style=nextline,itemsep=1pt,topsep=3pt]
  \item[CWE] CWE-829
  \item[OWASP] OWASP ASVS 5.0 --- dependency/supply-chain controls
  \item[CERT] Not mapped
  \item[CodeQL] \textbf{Requires another control}; Dependabot alerts/security updates + dependency review
  \item[Enforcement] Dependabot + dependency review + lockfiles + approved registries
  \item[Automation] Yes
  \item[Sources] \texttt{GH-DEP-REVIEW, GH-DEPENDABOT, NIST-SSDF, JAVA-SCG}
\end{description}
\end{rulebox}

\begin{rulebox}[title={\texttt{SCS-ALL-RESOURCE-001} \hfill \prioritybadge{High Priority}{High}}]
\textbf{Domain.} Resource Management

\textbf{Requirement.} Untrusted inputs MUST be bounded before allocation, decompression, regex evaluation, recursion, parsing, fan-out, or other resource-amplifying work; services SHOULD enforce timeouts and quotas.

\textbf{Rationale.} Small attacker requests can trigger disproportionate CPU, memory, storage, connection, or thread consumption.

\textbf{Risk.} Denial of service, algorithmic complexity abuse

\begin{description}[leftmargin=2.65cm,labelwidth=2.3cm,style=nextline,itemsep=1pt,topsep=3pt]
  \item[CWE] CWE-400, CWE-770, CWE-789, CWE-1333
  \item[OWASP] OWASP ASVS 5.0 --- availability/resource controls
  \item[CERT] Not mapped
  \item[CodeQL] \textbf{Partial}; Allocation-size, regex, decompression and other resource queries vary by language
  \item[Enforcement] CodeQL + fuzz/load tests + runtime quotas/timeouts
  \item[Automation] Partial
  \item[Sources] \texttt{JAVA-SCG, GO-SEC, NODE-SEC, MITRE-CWE}
\end{description}
\end{rulebox}

\begin{rulebox}[title={\texttt{SCS-ALL-CONC-001} \hfill \prioritybadge{High Priority}{High}}]
\textbf{Domain.} Concurrency / Parallelism

\textbf{Requirement.} Shared mutable state MUST use language-appropriate synchronization/ownership mechanisms; security decisions MUST NOT rely on filesystem or shared state remaining unchanged between separate check and use operations.

\textbf{Rationale.} Races can invalidate authorization, resource, memory, or filesystem assumptions.

\textbf{Risk.} Race conditions, TOCTOU, corrupted shared state

\begin{description}[leftmargin=2.65cm,labelwidth=2.3cm,style=nextline,itemsep=1pt,topsep=3pt]
  \item[CWE] CWE-362, CWE-367
  \item[OWASP] Context-dependent / not forced
  \item[CERT] SEI CERT concurrency rules where applicable
  \item[CodeQL] \textbf{Partial}; TOCTOU/static queries where available; dynamic race detectors for Go and native code
  \item[Enforcement] CodeQL + race detector/thread sanitizer where supported + concurrency tests
  \item[Automation] Partial
  \item[Sources] \texttt{GO-SEC, CERT-C, CERT-CPP, MITRE-CWE}
\end{description}
\end{rulebox}

\begin{rulebox}[title={\texttt{SCS-ALL-DYN-001} \hfill \prioritybadge{Critical Control}{Critical}}]
\textbf{Domain.} Dynamic Code Execution

\textbf{Requirement.} Untrusted data MUST NOT reach eval-like APIs, runtime code generation, dynamic loaders, reflective invocation targets, or interpreters unless constrained by a purpose-built parser and explicit allowlist.

\textbf{Rationale.} Dynamic execution turns attacker-controlled data into executable behavior.

\textbf{Risk.} Code injection, arbitrary execution

\begin{description}[leftmargin=2.65cm,labelwidth=2.3cm,style=nextline,itemsep=1pt,topsep=3pt]
  \item[CWE] CWE-94, CWE-95
  \item[OWASP] OWASP ASVS 5.0 --- injection
  \item[CERT] Not mapped
  \item[CodeQL] \textbf{Direct/Partial}; Code-injection / dynamic-code queries vary by language
  \item[Enforcement] CodeQL + code review + API restrictions
  \item[Automation] Partial
  \item[Sources] \texttt{GH-BUILTIN, MITRE-CWE, OWASP-ASVS}
\end{description}
\end{rulebox}

\begin{rulebox}[title={\texttt{SCS-ALL-FFI-001} \hfill \prioritybadge{Critical Control}{Critical}}]
\textbf{Domain.} Native / Foreign Interfaces

\textbf{Requirement.} FFI/native boundaries MUST be treated as unsafe trust boundaries: signatures, ownership, lifetime, nullability, buffer lengths, thread-safety, and error conventions MUST be validated and documented; a safe wrapper SHOULD isolate raw calls.

\textbf{Rationale.} Foreign runtimes can invalidate guarantees made by a memory-safe or type-safe language.

\textbf{Risk.} Memory corruption, lifetime violations, ABI mismatch

\begin{description}[leftmargin=2.65cm,labelwidth=2.3cm,style=nextline,itemsep=1pt,topsep=3pt]
  \item[CWE] CWE-119, CWE-416, CWE-704
  \item[OWASP] Context-dependent / not forced
  \item[CERT] SEI CERT / language-specific FFI guidance
  \item[CodeQL] \textbf{Partial}; Native-side CodeQL where source is analyzed; limited cross-language guarantees
  \item[Enforcement] Boundary review + native/static analysis + sanitizers + tests
  \item[Automation] Partial
  \item[Sources] \texttt{RUST-FFI, SWIFT-INTEROP, JAVA-SCG, DOTNET-SEC}
\end{description}
\end{rulebox}

\begin{rulebox}[title={\texttt{SCS-ALL-BUILD-001} \hfill \prioritybadge{High Priority}{High}}]
\textbf{Domain.} Compiler / Build Security

\textbf{Requirement.} Builds MUST use supported toolchains, fail on security-relevant static-analysis findings according to policy, preserve reproducible dependency state where the ecosystem supports lockfiles, and SHOULD run language-native security diagnostics/tests in CI.

\textbf{Rationale.} Security controls that are not reproducible and enforced in CI degrade over time and across developer environments.

\textbf{Risk.} Unreviewed defects, vulnerable dependencies, inconsistent builds

\begin{description}[leftmargin=2.65cm,labelwidth=2.3cm,style=nextline,itemsep=1pt,topsep=3pt]
  \item[CWE] CWE-1104 (where applicable)
  \item[OWASP] Context-dependent / not forced
  \item[CERT] Not mapped
  \item[CodeQL] \textbf{Requires another control}; CodeQL check runs plus language-native compiler/linter/test/dependency checks
  \item[Enforcement] CI/CD policy + CodeQL + ecosystem-native tooling
  \item[Automation] Yes/Partial
  \item[Sources] \texttt{NIST-SSDF, GH-LANG, GH-DEP-REVIEW, GO-SEC, TS-STRICT}
\end{description}
\end{rulebox}
\clearpage
\section{Language-Specific Standards}
Each language chapter combines the shared baseline with controls unique to that language, runtime, ecosystem, or interoperability model.
\clearpage
\subsection{C}
C provides direct memory access, pointer arithmetic, manual lifetime management, implementation/undefined behavior, variadic formatting, and low-level filesystem/process interfaces. Security therefore depends heavily on explicit bounds/lifetime/arithmetic invariants and toolchain verification. Primary authority: SEI CERT C Coding Standard (2016).
\begin{callout}{Teal}
\textbf{Authority classification:} Formal Secure Coding Standard\\
\textbf{Primary authority:} SEI CERT C Coding Standard (2016)\\
\textbf{Secondary authorities:} SEI CERT + NIST SSDF + MITRE CWE + OWASP where web-facing
\end{callout}
\subsubsection*{Language-specific secure coding requirements}

\begin{rulebox}[title={\texttt{SCS-C-MEM-001} \hfill \prioritybadge{Critical Control}{Critical}}]
\textbf{Domain.} Memory Safety

\textbf{Requirement.} Pointer arithmetic, array indexing, copy operations, and allocation sizes MUST be proven within valid object bounds before use.

\textbf{Rationale.} C permits unchecked pointer arithmetic and raw memory access.

\textbf{Risk.} Buffer overflow / out-of-bounds access

\begin{description}[leftmargin=2.65cm,labelwidth=2.3cm,style=nextline,itemsep=1pt,topsep=3pt]
  \item[CWE] CWE-119, CWE-787
  \item[OWASP] Context-dependent / not forced
  \item[CERT] SEI CERT ARR/MEM guidance
  \item[CodeQL] \textbf{Direct/Partial}; C/C++ bounds and buffer-overflow queries
  \item[Enforcement] CodeQL + compiler diagnostics + sanitizers + review
  \item[Automation] Partial
  \item[Sources] \texttt{CERT-C, GH-BUILTIN}
\end{description}
\end{rulebox}

\begin{rulebox}[title={\texttt{SCS-C-MEM-002} \hfill \prioritybadge{Critical Control}{Critical}}]
\textbf{Domain.} Memory Safety

\textbf{Requirement.} Allocated objects MUST have a single, auditable ownership/lifetime model; code MUST NOT dereference or free storage after its lifetime ends and MUST NOT free the same allocation twice.

\textbf{Rationale.} Manual lifetime management enables use-after-free and double-free defects.

\textbf{Risk.} Use-after-free / double free

\begin{description}[leftmargin=2.65cm,labelwidth=2.3cm,style=nextline,itemsep=1pt,topsep=3pt]
  \item[CWE] CWE-416, CWE-415
  \item[OWASP] Context-dependent / not forced
  \item[CERT] SEI CERT MEM guidance
  \item[CodeQL] \textbf{Direct}; cpp/use-after-free and related double-free queries
  \item[Enforcement] CodeQL + ASan-capable test builds + review
  \item[Automation] Yes/Partial
  \item[Sources] \texttt{CERT-C, GH-BUILTIN}
\end{description}
\end{rulebox}

\begin{rulebox}[title={\texttt{SCS-C-INPUT-001} \hfill \prioritybadge{High Priority}{High}}]
\textbf{Domain.} Input Validation

\textbf{Requirement.} Numeric sizes, offsets, indices, and conversions derived from untrusted input MUST be checked before arithmetic, narrowing conversion, allocation, indexing, or copy operations.

\textbf{Rationale.} Integer wrap/truncation can invalidate bounds checks and create undersized allocations.

\textbf{Risk.} Integer overflow/wrap/truncation

\begin{description}[leftmargin=2.65cm,labelwidth=2.3cm,style=nextline,itemsep=1pt,topsep=3pt]
  \item[CWE] CWE-190, CWE-681
  \item[OWASP] Context-dependent / not forced
  \item[CERT] SEI CERT INT guidance
  \item[CodeQL] \textbf{Partial}; C/C++ overflow and conversion queries
  \item[Enforcement] CodeQL + compiler diagnostics + boundary tests
  \item[Automation] Partial
  \item[Sources] \texttt{CERT-C, GH-BUILTIN}
\end{description}
\end{rulebox}

\begin{rulebox}[title={\texttt{SCS-C-OUTPUT-001} \hfill \prioritybadge{High Priority}{High}}]
\textbf{Domain.} Output Encoding

\textbf{Requirement.} Format strings passed to printf-family or equivalent variadic formatting functions MUST be compile-time constants or otherwise trusted; untrusted text MUST be passed as data arguments.

\textbf{Rationale.} Attacker-controlled format specifiers can read/write memory or disclose process state.

\textbf{Risk.} Uncontrolled format string

\begin{description}[leftmargin=2.65cm,labelwidth=2.3cm,style=nextline,itemsep=1pt,topsep=3pt]
  \item[CWE] CWE-134
  \item[OWASP] Context-dependent / not forced
  \item[CERT] SEI CERT FIO guidance
  \item[CodeQL] \textbf{Direct}; C/C++ uncontrolled-format-string query
  \item[Enforcement] CodeQL + compiler format warnings
  \item[Automation] Yes/Partial
  \item[Sources] \texttt{CERT-C, GH-BUILTIN}
\end{description}
\end{rulebox}

\begin{rulebox}[title={\texttt{SCS-C-FILE-001} \hfill \prioritybadge{High Priority}{High}}]
\textbf{Domain.} File / Path Security

\textbf{Requirement.} Security checks on filesystem objects MUST be coupled to the opened handle where possible; separate path checks followed by later path use MUST NOT be treated as stable authorization.

\textbf{Rationale.} Filesystem state can change between check and use.

\textbf{Risk.} TOCTOU race

\begin{description}[leftmargin=2.65cm,labelwidth=2.3cm,style=nextline,itemsep=1pt,topsep=3pt]
  \item[CWE] CWE-367
  \item[OWASP] Context-dependent / not forced
  \item[CERT] SEI CERT FIO/POS guidance
  \item[CodeQL] \textbf{Direct/Partial}; C/C++ TOCTOU queries
  \item[Enforcement] CodeQL + handle-based APIs + tests
  \item[Automation] Partial
  \item[Sources] \texttt{CERT-C, GH-BUILTIN}
\end{description}
\end{rulebox}

\begin{rulebox}[title={\texttt{SCS-C-BUILD-001} \hfill \prioritybadge{High Priority}{High}}]
\textbf{Domain.} Compiler / Build Security

\textbf{Requirement.} CI SHOULD include compiler diagnostics at a project-approved strict level plus memory/undefined-behavior instrumentation in dedicated test jobs where the selected toolchain supports it.

\textbf{Rationale.} Native-code defects often require compiler/runtime instrumentation beyond static analysis.

\textbf{Risk.} Latent memory/undefined-behavior defects

\begin{description}[leftmargin=2.65cm,labelwidth=2.3cm,style=nextline,itemsep=1pt,topsep=3pt]
  \item[CWE] Not forced / context-dependent
  \item[OWASP] Context-dependent / not forced
  \item[CERT] SEI CERT + toolchain documentation
  \item[CodeQL] \textbf{Requires another control}; ---
  \item[Enforcement] Compiler diagnostics + sanitizer test jobs + CodeQL
  \item[Automation] Yes/Partial
  \item[Sources] \texttt{CERT-C, NIST-SSDF}
\end{description}
\end{rulebox}
\subsubsection*{Representative implementation example}
\subsubsection*{Example --- bounds before copy}

\textbf{Noncompliant}\par
\begin{lstlisting}[language={C}]
void copy_name(char dst[32], const char *src, size_t n) {
    memcpy(dst, src, n);              /* n may exceed dst */
}
\end{lstlisting}

\textbf{Compliant}\par
\begin{lstlisting}[language={C}]
#include <string.h>
#include <stdbool.h>

bool copy_name(char dst[32], const char *src, size_t n) {
    if (n >= 32) return false;
    memcpy(dst, src, n);
    dst[n] = '\0';
    return true;
}
\end{lstlisting}

The trust boundary is the externally derived length \code{n}. The secure form establishes a destination bound before the copy. Real code must also validate that \code{src} is valid for at least \code{n} bytes. Relevant weakness: CWE-787/CWE-119. CodeQL can detect many out-of-bounds patterns, but dynamic sizes still require runtime checks and testing.\par
\subsubsection*{Recommended verification controls}
CodeQL C/C++ + strict compiler diagnostics + ASan/UBSan-capable test jobs where supported + fuzzing + dependency/secret controls + manual review of native boundaries.
\clearpage
\subsection{C++}
C++ adds RAII, templates, references, containers, smart pointers, exceptions, and a complex object/lifetime model on top of native memory access. Safer abstractions reduce risk but do not eliminate undefined behavior or dangling lifetime problems. Primary authority: SEI CERT C++ Coding Standard (2016).
\begin{callout}{Teal}
\textbf{Authority classification:} Formal Secure Coding Standard\\
\textbf{Primary authority:} SEI CERT C++ Coding Standard (2016)\\
\textbf{Secondary authorities:} SEI CERT + NIST SSDF + MITRE CWE + OWASP where web-facing
\end{callout}
\subsubsection*{Language-specific secure coding requirements}

\begin{rulebox}[title={\texttt{SCS-CPP-MEM-001} \hfill \prioritybadge{Critical Control}{Critical}}]
\textbf{Domain.} Memory Safety

\textbf{Requirement.} Owning resources SHOULD use RAII and standard ownership abstractions; raw pointers/references SHOULD be non-owning unless an explicitly reviewed invariant requires otherwise.

\textbf{Rationale.} RAII ties resource lifetime to object lifetime and reduces leaks, double deletion, and exceptional-path cleanup failures.

\textbf{Risk.} Use-after-free, double free, leaks

\begin{description}[leftmargin=2.65cm,labelwidth=2.3cm,style=nextline,itemsep=1pt,topsep=3pt]
  \item[CWE] CWE-416, CWE-415
  \item[OWASP] Context-dependent / not forced
  \item[CERT] SEI CERT MEM/ERR guidance
  \item[CodeQL] \textbf{Partial}; C/C++ lifetime, mismatched allocation/deallocation, and UAF queries
  \item[Enforcement] CodeQL + review + sanitizer tests
  \item[Automation] Partial
  \item[Sources] \texttt{CERT-CPP, GH-BUILTIN}
\end{description}
\end{rulebox}

\begin{rulebox}[title={\texttt{SCS-CPP-MEM-002} \hfill \prioritybadge{Critical Control}{Critical}}]
\textbf{Domain.} Memory Safety

\textbf{Requirement.} Code MUST respect object, reference, iterator, and view lifetimes; references/iterators/views MUST NOT escape beyond the lifetime or mutation rules of their backing object.

\textbf{Rationale.} C++ abstractions can preserve dangling references even when containers/smart pointers are used elsewhere.

\textbf{Risk.} Use-after-free / invalid iterator

\begin{description}[leftmargin=2.65cm,labelwidth=2.3cm,style=nextline,itemsep=1pt,topsep=3pt]
  \item[CWE] CWE-416, CWE-825
  \item[OWASP] Context-dependent / not forced
  \item[CERT] SEI CERT CTR/MEM guidance
  \item[CodeQL] \textbf{Direct/Partial}; cpp/use-after-free and iterator/lifetime queries
  \item[Enforcement] CodeQL + compiler/sanitizer + review
  \item[Automation] Partial
  \item[Sources] \texttt{CERT-CPP, GH-BUILTIN}
\end{description}
\end{rulebox}

\begin{rulebox}[title={\texttt{SCS-CPP-INPUT-001} \hfill \prioritybadge{High Priority}{High}}]
\textbf{Domain.} Input Validation

\textbf{Requirement.} Signed/unsigned conversions, narrowing conversions, container sizes, and arithmetic used for allocation/indexing MUST be range-checked before use.

\textbf{Rationale.} Implicit conversions can defeat bounds reasoning and produce wraparound.

\textbf{Risk.} Integer conversion/overflow

\begin{description}[leftmargin=2.65cm,labelwidth=2.3cm,style=nextline,itemsep=1pt,topsep=3pt]
  \item[CWE] CWE-190, CWE-681
  \item[OWASP] Context-dependent / not forced
  \item[CERT] SEI CERT INT guidance
  \item[CodeQL] \textbf{Partial}; C/C++ integer and allocation-size queries
  \item[Enforcement] CodeQL + compiler conversion warnings + tests
  \item[Automation] Partial
  \item[Sources] \texttt{CERT-CPP, GH-BUILTIN}
\end{description}
\end{rulebox}

\begin{rulebox}[title={\texttt{SCS-CPP-FFI-001} \hfill \prioritybadge{Critical Control}{Critical}}]
\textbf{Domain.} Native / Foreign Interfaces

\textbf{Requirement.} C APIs, variadic functions, raw allocation APIs, and ABI boundaries MUST be isolated behind typed wrappers that document ownership, lengths, nullability, and error behavior.

\textbf{Rationale.} Low-level interfaces bypass higher-level type/lifetime guarantees.

\textbf{Risk.} Memory corruption / ABI misuse

\begin{description}[leftmargin=2.65cm,labelwidth=2.3cm,style=nextline,itemsep=1pt,topsep=3pt]
  \item[CWE] CWE-119, CWE-704
  \item[OWASP] Context-dependent / not forced
  \item[CERT] SEI CERT API/MEM guidance
  \item[CodeQL] \textbf{Partial}; ---
  \item[Enforcement] Wrapper review + CodeQL + sanitizer tests
  \item[Automation] Partial
  \item[Sources] \texttt{CERT-CPP, NIST-SSDF}
\end{description}
\end{rulebox}

\begin{rulebox}[title={\texttt{SCS-CPP-CONC-001} \hfill \prioritybadge{High Priority}{High}}]
\textbf{Domain.} Concurrency / Parallelism

\textbf{Requirement.} Shared mutable state MUST be synchronized with standard concurrency primitives; lifetime and ownership MUST remain valid across asynchronous tasks and threads.

\textbf{Rationale.} Data races are undefined behavior in C++ and can corrupt security-relevant state.

\textbf{Risk.} Race condition

\begin{description}[leftmargin=2.65cm,labelwidth=2.3cm,style=nextline,itemsep=1pt,topsep=3pt]
  \item[CWE] CWE-362
  \item[OWASP] Context-dependent / not forced
  \item[CERT] SEI CERT CON guidance
  \item[CodeQL] \textbf{Partial}; ---
  \item[Enforcement] Thread sanitizer / tests + review; CodeQL where modeled
  \item[Automation] Partial
  \item[Sources] \texttt{CERT-CPP, MITRE-CWE}
\end{description}
\end{rulebox}

\begin{rulebox}[title={\texttt{SCS-CPP-BUILD-001} \hfill \prioritybadge{High Priority}{High}}]
\textbf{Domain.} Compiler / Build Security

\textbf{Requirement.} CI SHOULD combine CodeQL with strict compiler diagnostics and sanitizer/instrumented test configurations supported by the project toolchain.

\textbf{Rationale.} Template, conversion, lifetime, and memory defects need complementary analyses.

\textbf{Risk.} Latent native-code defects

\begin{description}[leftmargin=2.65cm,labelwidth=2.3cm,style=nextline,itemsep=1pt,topsep=3pt]
  \item[CWE] Not forced / context-dependent
  \item[OWASP] Context-dependent / not forced
  \item[CERT] SEI CERT + toolchain guidance
  \item[CodeQL] \textbf{Requires another control}; ---
  \item[Enforcement] Compiler + sanitizer tests + CodeQL
  \item[Automation] Yes/Partial
  \item[Sources] \texttt{CERT-CPP, NIST-SSDF}
\end{description}
\end{rulebox}
\subsubsection*{Representative implementation example}
\subsubsection*{Example --- explicit ownership}

\textbf{Noncompliant}\par
\begin{lstlisting}[language={C++}]
Widget* make_widget() { return new Widget(); }
// Callers must guess who owns/deletes the object.
\end{lstlisting}

\textbf{Compliant}\par
\begin{lstlisting}[language={C++}]
#include <memory>
std::unique_ptr<Widget> make_widget() {
    return std::make_unique<Widget>();
}
\end{lstlisting}

The secure form makes ownership explicit through RAII. It reduces exceptional-path leaks and ambiguous deletion responsibility; it does not by itself prevent all dangling non-owning references. Relevant weakness families include CWE-416/CWE-415.\par
\subsubsection*{Recommended verification controls}
CodeQL C/C++ + compiler warnings + ASan/UBSan/TSan-capable test jobs where supported + fuzzing + dependency/secret controls + ownership/lifetime review.
\clearpage
\subsection{C\#}
C\#/.NET normally provides managed memory and type safety, but unsafe code, P/Invoke, reflection, powerful serializers, web frameworks, and native interop create explicit security boundaries. Microsoft explicitly advises against BinaryFormatter and several legacy trust mechanisms. Primary authority: Microsoft .NET secure coding guidance.
\begin{callout}{Teal}
\textbf{Authority classification:} Vendor Secure Coding Guidelines\\
\textbf{Primary authority:} Microsoft .NET secure coding/security documentation\\
\textbf{Secondary authorities:} ASP.NET Core security docs + NIST/OWASP/MITRE + GitHub
\end{callout}
\subsubsection*{Language-specific secure coding requirements}

\begin{rulebox}[title={\texttt{SCS-CSHARP-DESER-001} \hfill \prioritybadge{Critical Control}{Critical}}]
\textbf{Domain.} Serialization / Deserialization

\textbf{Requirement.} Applications MUST NOT use BinaryFormatter for data processing and MUST NOT treat a binder or configuration switch as making BinaryFormatter deserialization safe.

\textbf{Rationale.} Microsoft states BinaryFormatter is insecure and cannot be made secure; attacker-controlled payloads can alter control flow and achieve RCE.

\textbf{Risk.} Unsafe deserialization

\begin{description}[leftmargin=2.65cm,labelwidth=2.3cm,style=nextline,itemsep=1pt,topsep=3pt]
  \item[CWE] CWE-502
  \item[OWASP] OWASP ASVS 5.0 --- deserialization
  \item[CERT] Not mapped
  \item[CodeQL] \textbf{Direct}; cs/unsafe-deserialization-untrusted-input
  \item[Enforcement] CodeQL + API ban/analyzer + migration to data-only serializers
  \item[Automation] Yes
  \item[Sources] \texttt{DOTNET-BF, GH-BUILTIN}
\end{description}
\end{rulebox}

\begin{rulebox}[title={\texttt{SCS-CSHARP-CRYPTO-001} \hfill \prioritybadge{Critical Control}{Critical}}]
\textbf{Domain.} Cryptography

\textbf{Requirement.} Security-sensitive random values MUST use RandomNumberGenerator or an appropriate higher-level .NET cryptographic API; System.Random MUST NOT be used for tokens, keys, salts, or nonces.

\textbf{Rationale.} General-purpose pseudo-random generators are predictable for adversarial security use.

\textbf{Risk.} Insufficiently random values

\begin{description}[leftmargin=2.65cm,labelwidth=2.3cm,style=nextline,itemsep=1pt,topsep=3pt]
  \item[CWE] CWE-338
  \item[OWASP] OWASP ASVS 5.0 --- cryptography
  \item[CERT] Not mapped
  \item[CodeQL] \textbf{Direct/Partial}; C\# insecure-randomness / weak-crypto queries
  \item[Enforcement] CodeQL + API restrictions + tests
  \item[Automation] Partial
  \item[Sources] \texttt{DOTNET-CRYPTO, GH-BUILTIN}
\end{description}
\end{rulebox}

\begin{rulebox}[title={\texttt{SCS-CSHARP-FFI-001} \hfill \prioritybadge{Critical Control}{Critical}}]
\textbf{Domain.} Native / Foreign Interfaces

\textbf{Requirement.} unsafe code and P/Invoke MUST be isolated, minimized, and reviewed as native trust boundaries with explicit buffer, lifetime, encoding, and ownership contracts.

\textbf{Rationale.} Managed runtime safety does not extend through unchecked pointer or native ABI operations.

\textbf{Risk.} Memory corruption / marshaling errors

\begin{description}[leftmargin=2.65cm,labelwidth=2.3cm,style=nextline,itemsep=1pt,topsep=3pt]
  \item[CWE] CWE-119, CWE-704
  \item[OWASP] Context-dependent / not forced
  \item[CERT] Not mapped
  \item[CodeQL] \textbf{Partial}; ---
  \item[Enforcement] Code review + native-side analysis + tests
  \item[Automation] Partial
  \item[Sources] \texttt{DOTNET-SEC, NIST-SSDF}
\end{description}
\end{rulebox}

\begin{rulebox}[title={\texttt{SCS-CSHARP-WEB-001} \hfill \prioritybadge{Critical Control}{Critical}}]
\textbf{Domain.} Web Application Security

\textbf{Requirement.} ASP.NET Core applications using cookie-authenticated state-changing requests MUST use the framework anti-forgery mechanisms where the request pattern is CSRF-relevant.

\textbf{Rationale.} Browsers automatically attach ambient credentials, enabling forged cross-site state-changing requests.

\textbf{Risk.} CSRF

\begin{description}[leftmargin=2.65cm,labelwidth=2.3cm,style=nextline,itemsep=1pt,topsep=3pt]
  \item[CWE] CWE-352
  \item[OWASP] OWASP ASVS 5.0 --- CSRF
  \item[CERT] Not mapped
  \item[CodeQL] \textbf{Partial}; ASP.NET CSRF-related CodeQL queries where modeled
  \item[Enforcement] Framework configuration + integration/DAST tests + CodeQL
  \item[Automation] Partial
  \item[Sources] \texttt{ASPNET-CSRF, OWASP-ASVS, GH-BUILTIN}
\end{description}
\end{rulebox}

\begin{rulebox}[title={\texttt{SCS-CSHARP-AUTHZ-001} \hfill \prioritybadge{Critical Control}{Critical}}]
\textbf{Domain.} Authorization

\textbf{Requirement.} ASP.NET authorization MUST be enforced with server-side policies/handlers at protected endpoints and resources; UI visibility MUST NOT be treated as authorization.

\textbf{Rationale.} Client-visible controls do not prevent direct requests to protected resources.

\textbf{Risk.} Missing authorization / IDOR

\begin{description}[leftmargin=2.65cm,labelwidth=2.3cm,style=nextline,itemsep=1pt,topsep=3pt]
  \item[CWE] CWE-862, CWE-639
  \item[OWASP] OWASP ASVS 5.0 --- access control
  \item[CERT] Not mapped
  \item[CodeQL] \textbf{Partial}; C\# access-control / IDOR queries where framework modeling applies
  \item[Enforcement] Authorization tests + CodeQL + review
  \item[Automation] Partial
  \item[Sources] \texttt{DOTNET-SEC, OWASP-ASVS, GH-BUILTIN}
\end{description}
\end{rulebox}

\begin{rulebox}[title={\texttt{SCS-CSHARP-DYN-001} \hfill \prioritybadge{Critical Control}{Critical}}]
\textbf{Domain.} Dynamic Code Execution

\textbf{Requirement.} Untrusted data MUST NOT determine dynamically loaded assemblies/types, reflection targets, expressions, or process command strings without strict allowlisting and structured APIs.

\textbf{Rationale.} Reflection/dynamic loading and command construction can convert data into executable behavior.

\textbf{Risk.} Code/command injection

\begin{description}[leftmargin=2.65cm,labelwidth=2.3cm,style=nextline,itemsep=1pt,topsep=3pt]
  \item[CWE] CWE-78, CWE-94
  \item[OWASP] OWASP ASVS 5.0 --- injection
  \item[CERT] Not mapped
  \item[CodeQL] \textbf{Direct/Partial}; cs/command-line-injection plus dynamic-code queries
  \item[Enforcement] CodeQL + API restrictions + review
  \item[Automation] Partial
  \item[Sources] \texttt{GH-BUILTIN, DOTNET-SEC}
\end{description}
\end{rulebox}
\subsubsection*{Representative implementation example}
\subsubsection*{Example --- do not deserialize arbitrary objects}

\textbf{Noncompliant}\par
\begin{lstlisting}[language={[Sharp]C}]
#pragma warning disable SYSLIB0011
var formatter = new System.Runtime.Serialization.Formatters.Binary.BinaryFormatter();
var value = formatter.Deserialize(stream);
#pragma warning restore SYSLIB0011
\end{lstlisting}

\textbf{Compliant}\par
\begin{lstlisting}[language={[Sharp]C}]
using System.Text.Json;

var value = await JsonSerializer.DeserializeAsync<MyDto>(stream);
if (value is null) throw new InvalidDataException("Invalid payload");
Validate(value);
\end{lstlisting}

Microsoft states that \code{BinaryFormatter} is insecure and cannot be made secure. The replacement still needs domain/schema validation. Relevant weakness: CWE-502. Representative CodeQL query: \code{cs/unsafe-deserialization-untrusted-input}.\par
\subsubsection*{Recommended verification controls}
CodeQL C\# + .NET analyzers/compiler + unit/integration tests + ASP.NET security tests where applicable + dependency review/Dependabot + secret scanning + native-boundary review.
\clearpage
\subsection{Go}
Go provides memory safety for ordinary code, garbage collection, explicit error returns, goroutines/channels, and a security-focused standard library. Security still depends on correct templating, process invocation, TLS verification, cryptographic randomness, dependency hygiene, and concurrency testing. Primary authority: Go Security Best Practices plus standard-library docs.
\begin{callout}{Teal}
\textbf{Authority classification:} Official Security Guidance\\
\textbf{Primary authority:} Go Security Best Practices and standard-library docs\\
\textbf{Secondary authorities:} NIST/OWASP/MITRE + GitHub
\end{callout}
\subsubsection*{Language-specific secure coding requirements}

\begin{rulebox}[title={\texttt{SCS-GO-OUTPUT-001} \hfill \prioritybadge{High Priority}{High}}]
\textbf{Domain.} Output Encoding

\textbf{Requirement.} HTML responses containing untrusted values MUST use html/template or an equivalently context-aware escaping mechanism; text/template MUST NOT be substituted for HTML-context escaping.

\textbf{Rationale.} html/template is designed for safe HTML output contexts; plain text templating does not provide equivalent contextual escaping.

\textbf{Risk.} XSS

\begin{description}[leftmargin=2.65cm,labelwidth=2.3cm,style=nextline,itemsep=1pt,topsep=3pt]
  \item[CWE] CWE-79
  \item[OWASP] OWASP ASVS 5.0 --- output encoding
  \item[CERT] Not mapped
  \item[CodeQL] \textbf{Direct/Partial}; Go XSS queries
  \item[Enforcement] CodeQL + template API policy + tests
  \item[Automation] Partial
  \item[Sources] \texttt{GO-TEMPLATE, GH-BUILTIN}
\end{description}
\end{rulebox}

\begin{rulebox}[title={\texttt{SCS-GO-INJ-001} \hfill \prioritybadge{Critical Control}{Critical}}]
\textbf{Domain.} Injection Prevention

\textbf{Requirement.} Processes MUST be invoked with explicit executable and argument boundaries; untrusted data MUST NOT be embedded in a shell command string.

\textbf{Rationale.} Shell parsing gives metacharacters control over command structure.

\textbf{Risk.} OS command injection

\begin{description}[leftmargin=2.65cm,labelwidth=2.3cm,style=nextline,itemsep=1pt,topsep=3pt]
  \item[CWE] CWE-78
  \item[OWASP] Context-dependent / not forced
  \item[CERT] Not mapped
  \item[CodeQL] \textbf{Direct}; Go command-injection query
  \item[Enforcement] CodeQL + os/exec API policy
  \item[Automation] Yes/Partial
  \item[Sources] \texttt{GO-EXEC, GH-BUILTIN}
\end{description}
\end{rulebox}

\begin{rulebox}[title={\texttt{SCS-GO-NET-001} \hfill \prioritybadge{Critical Control}{Critical}}]
\textbf{Domain.} Network Communications

\textbf{Requirement.} Production TLS clients MUST NOT set InsecureSkipVerify to disable certificate/hostname verification unless a narrowly reviewed custom verification design provides equivalent authentication.

\textbf{Rationale.} Disabling peer verification exposes TLS sessions to impersonation.

\textbf{Risk.} Improper certificate validation

\begin{description}[leftmargin=2.65cm,labelwidth=2.3cm,style=nextline,itemsep=1pt,topsep=3pt]
  \item[CWE] CWE-295
  \item[OWASP] OWASP ASVS 5.0 --- communications
  \item[CERT] Not mapped
  \item[CodeQL] \textbf{Direct}; go/disabled-certificate-check
  \item[Enforcement] CodeQL required check + integration tests
  \item[Automation] Yes
  \item[Sources] \texttt{GH-BUILTIN, GO-SEC}
\end{description}
\end{rulebox}

\begin{rulebox}[title={\texttt{SCS-GO-CRYPTO-001} \hfill \prioritybadge{Critical Control}{Critical}}]
\textbf{Domain.} Cryptography

\textbf{Requirement.} Security-sensitive random values MUST use crypto/rand or an approved higher-level cryptographic API, not math/rand.

\textbf{Rationale.} math/rand is not a cryptographically secure source.

\textbf{Risk.} Predictable tokens/keys

\begin{description}[leftmargin=2.65cm,labelwidth=2.3cm,style=nextline,itemsep=1pt,topsep=3pt]
  \item[CWE] CWE-338
  \item[OWASP] Context-dependent / not forced
  \item[CERT] Not mapped
  \item[CodeQL] \textbf{Direct/Partial}; Go insecure-randomness queries
  \item[Enforcement] CodeQL + API policy
  \item[Automation] Partial
  \item[Sources] \texttt{GO-RAND, GH-BUILTIN}
\end{description}
\end{rulebox}

\begin{rulebox}[title={\texttt{SCS-GO-CONC-001} \hfill \prioritybadge{High Priority}{High}}]
\textbf{Domain.} Concurrency / Parallelism

\textbf{Requirement.} CI MUST exercise concurrency-sensitive packages with the Go race detector where practical, and code MUST synchronize shared mutable state accessed by multiple goroutines.

\textbf{Rationale.} The race detector finds executed data races that static checks may miss.

\textbf{Risk.} Data race

\begin{description}[leftmargin=2.65cm,labelwidth=2.3cm,style=nextline,itemsep=1pt,topsep=3pt]
  \item[CWE] CWE-362
  \item[OWASP] Context-dependent / not forced
  \item[CERT] Not mapped
  \item[CodeQL] \textbf{Requires another control}; ---
  \item[Enforcement] go test -race + tests + CodeQL for other issues
  \item[Automation] Yes/Partial
  \item[Sources] \texttt{GO-SEC}
\end{description}
\end{rulebox}

\begin{rulebox}[title={\texttt{SCS-GO-DEP-001} \hfill \prioritybadge{High Priority}{High}}]
\textbf{Domain.} Dependency / Package Security

\textbf{Requirement.} CI SHOULD run govulncheck and MUST review reachable vulnerable dependencies; dependency upgrades SHOULD be tested and combined with GitHub dependency review/Dependabot where enabled.

\textbf{Rationale.} govulncheck uses the Go vulnerability database and can identify vulnerabilities affecting source/binaries.

\textbf{Risk.} Known-vulnerable dependency

\begin{description}[leftmargin=2.65cm,labelwidth=2.3cm,style=nextline,itemsep=1pt,topsep=3pt]
  \item[CWE] CWE-1104
  \item[OWASP] Context-dependent / not forced
  \item[CERT] Not mapped
  \item[CodeQL] \textbf{Requires another control}; ---
  \item[Enforcement] govulncheck + Dependabot + dependency review
  \item[Automation] Yes
  \item[Sources] \texttt{GO-SEC, GH-DEP-REVIEW, GH-DEPENDABOT}
\end{description}
\end{rulebox}
\subsubsection*{Representative implementation example}
\subsubsection*{Example --- contextual HTML escaping}

\textbf{Noncompliant}\par
\begin{lstlisting}
t := templateText.Must(templateText.New("x").Parse(`<p>{{.}}</p>`))
_ = t.Execute(w, userInput)
\end{lstlisting}

\textbf{Compliant}\par
\begin{lstlisting}
t := templateHTML.Must(templateHTML.New("x").Parse(`<p>{{.}}</p>`))
_ = t.Execute(w, userInput)
\end{lstlisting}

where the imports are aliased:\par
\begin{lstlisting}
import (
    templateHTML "html/template"
    templateText "text/template"
)
\end{lstlisting}

For HTML output, \code{html/template} performs contextual escaping designed for HTML. The trust boundary is \code{userInput}. Relevant weakness: CWE-79.\par
\subsubsection*{Recommended verification controls}
CodeQL Go + \code{go vet} + \code{go test -race} + fuzzing + \code{govulncheck} + dependency review/Dependabot + secret scanning.
\clearpage
\subsection{Java}
Java provides type safety, automatic memory management, and bounds checking, but Oracle stresses that secure coding remains necessary. Major risk areas include trust-boundary validation, resource exhaustion, deserialization, injection, XML/JNDI/reflection, access control, third-party libraries, and JNI. Primary authorities: SEI CERT Oracle Coding Standard for Java, augmented by Oracle Secure Coding Guidelines for current Java SE guidance.
\begin{callout}{Teal}
\textbf{Authority classification:} Formal Secure Coding Standard\\
\textbf{Primary authority:} SEI CERT Oracle Coding Standard for Java + Oracle Secure Coding Guidelines for Java SE\\
\textbf{Secondary authorities:} Oracle current Java SE guidance + NIST/OWASP/MITRE + GitHub
\end{callout}
\subsubsection*{Language-specific secure coding requirements}

\begin{rulebox}[title={\texttt{SCS-JAVA-DESER-001} \hfill \prioritybadge{Critical Control}{Critical}}]
\textbf{Domain.} Serialization / Deserialization

\textbf{Requirement.} Untrusted input MUST NOT be passed to Java native/object deserialization or other arbitrary-object deserializers; use constrained data formats and explicit type/schema handling.

\textbf{Rationale.} Deserialization can instantiate attacker-influenced object graphs and trigger gadget behavior.

\textbf{Risk.} Unsafe deserialization

\begin{description}[leftmargin=2.65cm,labelwidth=2.3cm,style=nextline,itemsep=1pt,topsep=3pt]
  \item[CWE] CWE-502
  \item[OWASP] OWASP ASVS 5.0 --- deserialization
  \item[CERT] SEI CERT Java serialization guidance
  \item[CodeQL] \textbf{Direct}; java/unsafe-deserialization
  \item[Enforcement] CodeQL + serializer allowlist/schema + review
  \item[Automation] Yes
  \item[Sources] \texttt{JAVA-SCG, GH-BUILTIN}
\end{description}
\end{rulebox}

\begin{rulebox}[title={\texttt{SCS-JAVA-INJ-001} \hfill \prioritybadge{Critical Control}{Critical}}]
\textbf{Domain.} Injection Prevention

\textbf{Requirement.} SQL MUST use PreparedStatement or equivalent parameter binding; process execution MUST separate executable/arguments and avoid shell interpretation of untrusted input.

\textbf{Rationale.} Structured APIs keep data from becoming interpreter syntax.

\textbf{Risk.} SQL / command injection

\begin{description}[leftmargin=2.65cm,labelwidth=2.3cm,style=nextline,itemsep=1pt,topsep=3pt]
  \item[CWE] CWE-89, CWE-78
  \item[OWASP] OWASP ASVS 5.0 --- injection
  \item[CERT] Not mapped
  \item[CodeQL] \textbf{Direct}; java/sql-injection; java/command-line-injection
  \item[Enforcement] CodeQL + API policy
  \item[Automation] Yes
  \item[Sources] \texttt{JAVA-SCG, GH-BUILTIN}
\end{description}
\end{rulebox}

\begin{rulebox}[title={\texttt{SCS-JAVA-INPUT-001} \hfill \prioritybadge{High Priority}{High}}]
\textbf{Domain.} Input Validation

\textbf{Requirement.} Inputs crossing trust boundaries MUST be validated before reflection, class/resource lookup, filesystem, XML, regex, allocation, or security-sensitive API use.

\textbf{Rationale.} Java memory/type safety does not establish that external data is safe for downstream interpreters or resources.

\textbf{Risk.} Multiple injection/resource weaknesses

\begin{description}[leftmargin=2.65cm,labelwidth=2.3cm,style=nextline,itemsep=1pt,topsep=3pt]
  \item[CWE] CWE-20
  \item[OWASP] Context-dependent / not forced
  \item[CERT] Not mapped
  \item[CodeQL] \textbf{Partial}; Multiple Java taint/data-flow queries
  \item[Enforcement] CodeQL + parser/schema tests
  \item[Automation] Partial
  \item[Sources] \texttt{JAVA-SCG, GH-BUILTIN}
\end{description}
\end{rulebox}

\begin{rulebox}[title={\texttt{SCS-JAVA-RESOURCE-001} \hfill \prioritybadge{High Priority}{High}}]
\textbf{Domain.} Resource Management

\textbf{Requirement.} Untrusted sizes, compressed payloads, regexes, XML, object graphs, and other potentially amplifying inputs MUST have application-specific resource limits.

\textbf{Rationale.} Oracle specifically calls out decompression bombs, regex backtracking, XML expansion, deserialization graphs, and oversized dimensions as DoS vectors.

\textbf{Risk.} Resource exhaustion

\begin{description}[leftmargin=2.65cm,labelwidth=2.3cm,style=nextline,itemsep=1pt,topsep=3pt]
  \item[CWE] CWE-400, CWE-770, CWE-1333
  \item[OWASP] Context-dependent / not forced
  \item[CERT] Not mapped
  \item[CodeQL] \textbf{Partial}; Java resource/regex/XML queries
  \item[Enforcement] CodeQL + limits + fuzz/load tests
  \item[Automation] Partial
  \item[Sources] \texttt{JAVA-SCG, GH-BUILTIN}
\end{description}
\end{rulebox}

\begin{rulebox}[title={\texttt{SCS-JAVA-FFI-001} \hfill \prioritybadge{Critical Control}{Critical}}]
\textbf{Domain.} Native / Foreign Interfaces

\textbf{Requirement.} JNI/native integrations MUST validate all lengths, nullability, ownership, thread, exception, and lifetime assumptions at the boundary.

\textbf{Rationale.} JNI bypasses Java's normal memory/type safety properties.

\textbf{Risk.} Memory corruption / native misuse

\begin{description}[leftmargin=2.65cm,labelwidth=2.3cm,style=nextline,itemsep=1pt,topsep=3pt]
  \item[CWE] CWE-119, CWE-704
  \item[OWASP] Context-dependent / not forced
  \item[CERT] Not mapped
  \item[CodeQL] \textbf{Partial}; ---
  \item[Enforcement] Java + native CodeQL where both sources are available; boundary review
  \item[Automation] Partial
  \item[Sources] \texttt{JAVA-SCG}
\end{description}
\end{rulebox}

\begin{rulebox}[title={\texttt{SCS-JAVA-DEP-001} \hfill \prioritybadge{High Priority}{High}}]
\textbf{Domain.} Dependency / Package Security

\textbf{Requirement.} Third-party libraries/frameworks MUST be inventoried, security updates tracked, and vulnerable versions remediated or formally mitigated.

\textbf{Rationale.} Oracle explicitly requires tracking and timely application of third-party security updates.

\textbf{Risk.} Known-vulnerable dependency

\begin{description}[leftmargin=2.65cm,labelwidth=2.3cm,style=nextline,itemsep=1pt,topsep=3pt]
  \item[CWE] CWE-1104
  \item[OWASP] Context-dependent / not forced
  \item[CERT] Not mapped
  \item[CodeQL] \textbf{Requires another control}; ---
  \item[Enforcement] Dependency review + Dependabot + ecosystem SBOM/scanner
  \item[Automation] Yes
  \item[Sources] \texttt{JAVA-SCG, GH-DEP-REVIEW, GH-DEPENDABOT}
\end{description}
\end{rulebox}
\subsubsection*{Representative implementation example}
\subsubsection*{Example --- parameterized SQL}

\textbf{Noncompliant}\par
\begin{lstlisting}[language={Java}]
String sql = "SELECT * FROM accounts WHERE name = '" + name + "'";
try (Statement s = connection.createStatement()) {
    s.executeQuery(sql);
}
\end{lstlisting}

\textbf{Compliant}\par
\begin{lstlisting}[language={Java}]
String sql = "SELECT * FROM accounts WHERE name = ?";
try (PreparedStatement s = connection.prepareStatement(sql)) {
    s.setString(1, name);
    s.executeQuery();
}
\end{lstlisting}

\code{name} is untrusted data. Binding it as a parameter keeps it out of SQL syntax. Relevant weakness: CWE-89. Representative CodeQL query: \code{java/sql-injection}.\par
\subsubsection*{Recommended verification controls}
CodeQL Java/Kotlin + compiler/build checks + unit/integration/fuzz tests + dependency review/Dependabot + secret scanning + framework/serialization/JNI review.
\clearpage
\subsection{Kotlin}
Kotlin improves null-safety and expressiveness but JVM applications inherit Java runtime/library risks. Platform types, \code{!!}, Java interoperability, reflection, serializers, and native/JNI code can escape Kotlin's normal guarantees. Primary authority: derived from Kotlin official security/null-safety docs plus Java/JVM guidance.
\begin{callout}{Teal}
\textbf{Authority classification:} Derived Secure Coding Baseline\\
\textbf{Primary authority:} Kotlin security/null-safety docs + JVM guidance\\
\textbf{Secondary authorities:} Oracle Java guidance + NIST/OWASP/MITRE + GitHub
\end{callout}
\subsubsection*{Language-specific secure coding requirements}

\begin{rulebox}[title={\texttt{SCS-KOTLIN-INPUT-001} \hfill \prioritybadge{High Priority}{High}}]
\textbf{Domain.} Input Validation

\textbf{Requirement.} External data MUST be validated at runtime; Kotlin type declarations and casts MUST NOT be treated as validation of network, file, database, or deserialized input.

\textbf{Rationale.} Type information is not a runtime trust decision for external data.

\textbf{Risk.} Improper input validation

\begin{description}[leftmargin=2.65cm,labelwidth=2.3cm,style=nextline,itemsep=1pt,topsep=3pt]
  \item[CWE] CWE-20
  \item[OWASP] OWASP ASVS 5.0 --- validation
  \item[CERT] Not mapped
  \item[CodeQL] \textbf{Partial}; Java/Kotlin data-flow queries where modeled
  \item[Enforcement] Schema/domain validation + CodeQL
  \item[Automation] Partial
  \item[Sources] \texttt{KOTLIN-SEC, OWASP-ASVS, GH-BUILTIN}
\end{description}
\end{rulebox}

\begin{rulebox}[title={\texttt{SCS-KOTLIN-MEM-001} \hfill \prioritybadge{High Priority}{High}}]
\textbf{Domain.} Memory Safety

\textbf{Requirement.} The not-null assertion operator (!!) SHOULD NOT be used on untrusted or platform-typed values unless a prior invariant check proves non-nullness.

\textbf{Rationale.} Kotlin null safety can be bypassed by explicit assertions and Java interoperability/platform types.

\textbf{Risk.} Null dereference / availability issue

\begin{description}[leftmargin=2.65cm,labelwidth=2.3cm,style=nextline,itemsep=1pt,topsep=3pt]
  \item[CWE] CWE-476
  \item[OWASP] Context-dependent / not forced
  \item[CERT] Not mapped
  \item[CodeQL] \textbf{Partial}; Java/Kotlin nullness-related queries where applicable
  \item[Enforcement] Compiler + review + tests
  \item[Automation] Partial
  \item[Sources] \texttt{KOTLIN-NULL, GH-BUILTIN}
\end{description}
\end{rulebox}

\begin{rulebox}[title={\texttt{SCS-KOTLIN-FFI-001} \hfill \prioritybadge{Critical Control}{Critical}}]
\textbf{Domain.} Native / Foreign Interfaces

\textbf{Requirement.} Kotlin/JVM code calling Java/native APIs MUST adopt the callee's security contract, including nullability, serialization, reflection, JNI, and exception behavior.

\textbf{Rationale.} Kotlin's source-level guarantees do not automatically constrain Java/native implementations.

\textbf{Risk.} Interop contract failure

\begin{description}[leftmargin=2.65cm,labelwidth=2.3cm,style=nextline,itemsep=1pt,topsep=3pt]
  \item[CWE] CWE-704, CWE-119
  \item[OWASP] Context-dependent / not forced
  \item[CERT] Not mapped
  \item[CodeQL] \textbf{Partial}; ---
  \item[Enforcement] Java/Kotlin CodeQL + boundary review
  \item[Automation] Partial
  \item[Sources] \texttt{KOTLIN-NULL, JAVA-SCG}
\end{description}
\end{rulebox}

\begin{rulebox}[title={\texttt{SCS-KOTLIN-DEP-001} \hfill \prioritybadge{High Priority}{High}}]
\textbf{Domain.} Dependency / Package Security

\textbf{Requirement.} Projects MUST track supported Kotlin release lines and current dependency security status; unsupported or vulnerable runtime/library versions MUST be upgraded or formally mitigated.

\textbf{Rationale.} JetBrains recommends current Kotlin releases and dependencies and documents a JVM stdlib support window beginning with Kotlin 2.4.0.

\textbf{Risk.} Known-vulnerable dependency

\begin{description}[leftmargin=2.65cm,labelwidth=2.3cm,style=nextline,itemsep=1pt,topsep=3pt]
  \item[CWE] CWE-1104
  \item[OWASP] Context-dependent / not forced
  \item[CERT] Not mapped
  \item[CodeQL] \textbf{Requires another control}; ---
  \item[Enforcement] Dependabot/dependency review + ecosystem vulnerability scanning
  \item[Automation] Yes
  \item[Sources] \texttt{KOTLIN-SEC, GH-DEP-REVIEW, GH-DEPENDABOT}
\end{description}
\end{rulebox}

\begin{rulebox}[title={\texttt{SCS-KOTLIN-DESER-001} \hfill \prioritybadge{Critical Control}{Critical}}]
\textbf{Domain.} Serialization / Deserialization

\textbf{Requirement.} JVM deserialization and third-party serialization libraries used from Kotlin MUST follow the same untrusted-input restrictions as Java; arbitrary-object deserialization MUST NOT accept untrusted payloads.

\textbf{Rationale.} Kotlin/JVM shares the runtime and library attack surface of Java.

\textbf{Risk.} Unsafe deserialization

\begin{description}[leftmargin=2.65cm,labelwidth=2.3cm,style=nextline,itemsep=1pt,topsep=3pt]
  \item[CWE] CWE-502
  \item[OWASP] OWASP ASVS 5.0 --- deserialization
  \item[CERT] Not mapped
  \item[CodeQL] \textbf{Direct/Partial}; java/unsafe-deserialization where the modeled JVM API is used
  \item[Enforcement] CodeQL + serializer policy
  \item[Automation] Partial
  \item[Sources] \texttt{JAVA-SCG, GH-BUILTIN}
\end{description}
\end{rulebox}
\subsubsection*{Representative implementation example}
\subsubsection*{Example --- parse at the trust boundary}

\textbf{Noncompliant}\par
\begin{lstlisting}
val count: Int = request.queryParameters["count"]!!.toInt()
\end{lstlisting}

\textbf{Compliant}\par
\begin{lstlisting}
val count = request.queryParameters["count"]
    ?.toIntOrNull()
    ?.takeIf { it in 1..100 }
    ?: return badRequest()
\end{lstlisting}

The secure form treats missing, malformed, and out-of-range input as expected boundary conditions rather than relying on \code{!!} and unchecked conversion. This is runtime validation, not merely type annotation.\par
\subsubsection*{Recommended verification controls}
CodeQL Java/Kotlin + Kotlin compiler strictness/nullability discipline + JVM dependency scanning + tests + dependency review/Dependabot + secret scanning + Java/native-boundary review.
\clearpage
\subsection{JavaScript}
JavaScript is dynamically typed and highly dependent on runtime context. Browser DOM APIs, server runtimes, prototypes, dynamic code execution, shell/process APIs, templating, async resource behavior, and npm supply-chain semantics dominate the security model. Primary authority: derived from Node.js official security/API docs plus web standards guidance.
\begin{callout}{Teal}
\textbf{Authority classification:} Derived Secure Coding Baseline\\
\textbf{Primary authority:} Node.js security/API docs + ECMAScript/runtime semantics\\
\textbf{Secondary authorities:} OWASP ASVS + NIST/MITRE + GitHub; browser framework docs as applicable
\end{callout}
\subsubsection*{Language-specific secure coding requirements}

\begin{rulebox}[title={\texttt{SCS-JS-DYN-001} \hfill \prioritybadge{Critical Control}{Critical}}]
\textbf{Domain.} Dynamic Code Execution

\textbf{Requirement.} Untrusted strings MUST NOT be passed to eval, Function constructors, dynamic code loaders, template compilation, or equivalent executable sinks.

\textbf{Rationale.} JavaScript's dynamic execution features can interpret attacker input directly as code.

\textbf{Risk.} Code injection

\begin{description}[leftmargin=2.65cm,labelwidth=2.3cm,style=nextline,itemsep=1pt,topsep=3pt]
  \item[CWE] CWE-94, CWE-95
  \item[OWASP] OWASP ASVS 5.0 --- injection
  \item[CERT] Not mapped
  \item[CodeQL] \textbf{Direct}; JavaScript code-injection queries
  \item[Enforcement] CodeQL + lint/API restrictions
  \item[Automation] Yes/Partial
  \item[Sources] \texttt{NODE-SEC, GH-BUILTIN}
\end{description}
\end{rulebox}

\begin{rulebox}[title={\texttt{SCS-JS-INPUT-001} \hfill \prioritybadge{High Priority}{High}}]
\textbf{Domain.} Input Validation

\textbf{Requirement.} When copying or merging untrusted object properties, dangerous prototype-related keys MUST be rejected or the application MUST use data structures/APIs that do not mutate object prototypes.

\textbf{Rationale.} Prototype pollution can change application behavior globally and create downstream injection or authorization failures.

\textbf{Risk.} Prototype pollution

\begin{description}[leftmargin=2.65cm,labelwidth=2.3cm,style=nextline,itemsep=1pt,topsep=3pt]
  \item[CWE] CWE-1321, CWE-915
  \item[OWASP] Context-dependent / not forced
  \item[CERT] Not mapped
  \item[CodeQL] \textbf{Direct}; js/prototype-polluting-assignment
  \item[Enforcement] CodeQL + safe merge/object patterns + tests
  \item[Automation] Yes
  \item[Sources] \texttt{NODE-SEC, GH-BUILTIN}
\end{description}
\end{rulebox}

\begin{rulebox}[title={\texttt{SCS-JS-INJ-001} \hfill \prioritybadge{Critical Control}{Critical}}]
\textbf{Domain.} Injection Prevention

\textbf{Requirement.} Node.js child processes MUST use executable/argument APIs without shell interpretation for untrusted arguments; shell strings MUST NOT contain untrusted data.

\textbf{Rationale.} Shell metacharacters can transform input into commands.

\textbf{Risk.} Command injection

\begin{description}[leftmargin=2.65cm,labelwidth=2.3cm,style=nextline,itemsep=1pt,topsep=3pt]
  \item[CWE] CWE-78
  \item[OWASP] Context-dependent / not forced
  \item[CERT] Not mapped
  \item[CodeQL] \textbf{Direct}; JavaScript command-injection queries
  \item[Enforcement] CodeQL + child\_process API policy
  \item[Automation] Yes/Partial
  \item[Sources] \texttt{NODE-CHILD, GH-BUILTIN}
\end{description}
\end{rulebox}

\begin{rulebox}[title={\texttt{SCS-JS-WEB-001} \hfill \prioritybadge{Critical Control}{Critical}}]
\textbf{Domain.} Web Application Security

\textbf{Requirement.} Browser/server JavaScript MUST keep untrusted data out of unsafe HTML/DOM/script sinks and use framework/context-aware escaping APIs.

\textbf{Rationale.} DOM and server-template sinks can interpret data as markup or code.

\textbf{Risk.} XSS / DOM XSS

\begin{description}[leftmargin=2.65cm,labelwidth=2.3cm,style=nextline,itemsep=1pt,topsep=3pt]
  \item[CWE] CWE-79
  \item[OWASP] OWASP ASVS 5.0 --- output encoding
  \item[CERT] Not mapped
  \item[CodeQL] \textbf{Direct}; JavaScript XSS / unsafe DOM queries
  \item[Enforcement] CodeQL + framework-safe rendering + DAST
  \item[Automation] Yes/Partial
  \item[Sources] \texttt{OWASP-ASVS, GH-BUILTIN}
\end{description}
\end{rulebox}

\begin{rulebox}[title={\texttt{SCS-JS-RESOURCE-001} \hfill \prioritybadge{High Priority}{High}}]
\textbf{Domain.} Resource Management

\textbf{Requirement.} Node.js network services MUST bound request body sizes, timeouts, socket/connection use, regex complexity, and other attacker-controlled resource costs.

\textbf{Rationale.} The event loop and shared process resources can be monopolized by adversarial work.

\textbf{Risk.} Denial of service

\begin{description}[leftmargin=2.65cm,labelwidth=2.3cm,style=nextline,itemsep=1pt,topsep=3pt]
  \item[CWE] CWE-400, CWE-1333
  \item[OWASP] Context-dependent / not forced
  \item[CERT] Not mapped
  \item[CodeQL] \textbf{Partial}; JavaScript regex/resource queries
  \item[Enforcement] CodeQL + runtime limits + load/fuzz tests
  \item[Automation] Partial
  \item[Sources] \texttt{NODE-SEC, GH-BUILTIN}
\end{description}
\end{rulebox}

\begin{rulebox}[title={\texttt{SCS-JS-DEP-001} \hfill \prioritybadge{Critical Control}{Critical}}]
\textbf{Domain.} Dependency / Package Security

\textbf{Requirement.} Package installation in CI MUST use committed lockfiles and controlled registries; package lifecycle scripts and dependency updates SHOULD be reviewed according to project supply-chain policy.

\textbf{Rationale.} The npm ecosystem executes third-party code during build/install and has a large transitive dependency surface.

\textbf{Risk.} Supply-chain compromise

\begin{description}[leftmargin=2.65cm,labelwidth=2.3cm,style=nextline,itemsep=1pt,topsep=3pt]
  \item[CWE] CWE-829
  \item[OWASP] Context-dependent / not forced
  \item[CERT] Not mapped
  \item[CodeQL] \textbf{Requires another control}; ---
  \item[Enforcement] Dependabot + dependency review + lockfile-controlled CI
  \item[Automation] Yes
  \item[Sources] \texttt{NODE-SEC, GH-DEP-REVIEW, GH-DEPENDABOT}
\end{description}
\end{rulebox}
\subsubsection*{Representative implementation example}
\subsubsection*{Example --- avoid shell interpretation}

\textbf{Noncompliant}\par
\begin{lstlisting}[language={Java}]
exec(`git show ${userRef}`, callback);
\end{lstlisting}

\textbf{Compliant}\par
\begin{lstlisting}[language={Java}]
execFile("git", ["show", userRef], callback);
\end{lstlisting}

The secure form preserves executable/argument boundaries instead of asking a shell to parse an attacker-influenced command string. The application still needs business validation of which refs are permitted. Relevant weakness: CWE-78.\par
\subsubsection*{Recommended verification controls}
CodeQL JavaScript/TypeScript + runtime/framework linting/tests + DAST/IAST for web apps + dependency review/Dependabot + secret scanning + lockfile/install-policy enforcement.
\clearpage
\subsection{TypeScript}
TypeScript strengthens compile-time checking but emits JavaScript. \code{any} disables checking; assertions are erased; external data remains unvalidated until runtime checks occur. All JavaScript runtime security controls continue to apply. Primary authority: TypeScript Handbook/TSConfig plus the JavaScript runtime baseline.
\begin{callout}{Teal}
\textbf{Authority classification:} Derived Secure Coding Baseline\\
\textbf{Primary authority:} TypeScript Handbook + TSConfig + JavaScript runtime baseline\\
\textbf{Secondary authorities:} Node/browser guidance + OWASP/NIST/MITRE + GitHub
\end{callout}
\subsubsection*{Language-specific secure coding requirements}

\begin{rulebox}[title={\texttt{SCS-TS-INPUT-001} \hfill \prioritybadge{High Priority}{High}}]
\textbf{Domain.} Input Validation

\textbf{Requirement.} Values from untrusted runtime boundaries SHOULD enter as unknown (or an equivalent schema-validation boundary), be validated/narrowed, and only then be assigned domain types.

\textbf{Rationale.} unknown forces narrowing while any disables type checking; neither a TypeScript annotation nor interface validates wire data.

\textbf{Risk.} Improper input validation

\begin{description}[leftmargin=2.65cm,labelwidth=2.3cm,style=nextline,itemsep=1pt,topsep=3pt]
  \item[CWE] CWE-20
  \item[OWASP] OWASP ASVS 5.0 --- validation
  \item[CERT] Not mapped
  \item[CodeQL] \textbf{Partial}; JavaScript/TypeScript data-flow queries
  \item[Enforcement] Runtime schema/domain validation + compiler + CodeQL
  \item[Automation] Partial
  \item[Sources] \texttt{TS-TYPES, OWASP-ASVS}
\end{description}
\end{rulebox}

\begin{rulebox}[title={\texttt{SCS-TS-INPUT-002} \hfill \prioritybadge{High Priority}{High}}]
\textbf{Domain.} Input Validation

\textbf{Requirement.} Type assertions (\code{as T}) and non-null assertions (\code{!}) MUST NOT be used as substitutes for runtime validation of attacker-controlled values.

\textbf{Rationale.} Type assertions are erased and perform no runtime check.

\textbf{Risk.} Improper validation / runtime type confusion

\begin{description}[leftmargin=2.65cm,labelwidth=2.3cm,style=nextline,itemsep=1pt,topsep=3pt]
  \item[CWE] CWE-20, CWE-704
  \item[OWASP] Context-dependent / not forced
  \item[CERT] Not mapped
  \item[CodeQL] \textbf{Partial}; ---
  \item[Enforcement] Compiler/review + runtime validation tests
  \item[Automation] Partial
  \item[Sources] \texttt{TS-TYPES}
\end{description}
\end{rulebox}

\begin{rulebox}[title={\texttt{SCS-TS-BUILD-001} \hfill \prioritybadge{High Priority}{High}}]
\textbf{Domain.} Compiler / Build Security

\textbf{Requirement.} Security-sensitive TypeScript projects SHOULD enable \code{strict}; exceptions to strict-family checks SHOULD be explicit and reviewed.

\textbf{Rationale.} Strict mode enables stronger compile-time correctness guarantees including nullability and implicit-any checks.

\textbf{Risk.} Latent type/nullability defects

\begin{description}[leftmargin=2.65cm,labelwidth=2.3cm,style=nextline,itemsep=1pt,topsep=3pt]
  \item[CWE] Not forced / context-dependent
  \item[OWASP] Context-dependent / not forced
  \item[CERT] Not mapped
  \item[CodeQL] \textbf{Requires another control}; ---
  \item[Enforcement] tsc --noEmit/CI + CodeQL
  \item[Automation] Yes
  \item[Sources] \texttt{TS-STRICT}
\end{description}
\end{rulebox}

\begin{rulebox}[title={\texttt{SCS-TS-DYN-001} \hfill \prioritybadge{Critical Control}{Critical}}]
\textbf{Domain.} Dynamic Code Execution

\textbf{Requirement.} Type safety MUST NOT be assumed to mitigate JavaScript runtime injection; TypeScript code MUST follow the JavaScript dynamic-execution, DOM, command, and prototype-pollution controls.

\textbf{Rationale.} TypeScript emits JavaScript and its type annotations do not create runtime security checks.

\textbf{Risk.} Code/DOM/command injection

\begin{description}[leftmargin=2.65cm,labelwidth=2.3cm,style=nextline,itemsep=1pt,topsep=3pt]
  \item[CWE] CWE-79, CWE-78, CWE-94
  \item[OWASP] Context-dependent / not forced
  \item[CERT] Not mapped
  \item[CodeQL] \textbf{Direct/Partial}; JavaScript/TypeScript injection and prototype-pollution queries
  \item[Enforcement] CodeQL + JavaScript runtime controls
  \item[Automation] Partial
  \item[Sources] \texttt{TS-TYPES, GH-BUILTIN}
\end{description}
\end{rulebox}

\begin{rulebox}[title={\texttt{SCS-TS-DEP-001} \hfill \prioritybadge{Critical Control}{Critical}}]
\textbf{Domain.} Dependency / Package Security

\textbf{Requirement.} Type declarations MUST NOT be treated as evidence that a package is trustworthy; packages, lockfiles, registries, install scripts, and transitive dependencies MUST follow the JavaScript supply-chain controls.

\textbf{Rationale.} Types describe interfaces, not package provenance or runtime integrity.

\textbf{Risk.} Supply-chain compromise

\begin{description}[leftmargin=2.65cm,labelwidth=2.3cm,style=nextline,itemsep=1pt,topsep=3pt]
  \item[CWE] CWE-829
  \item[OWASP] Context-dependent / not forced
  \item[CERT] Not mapped
  \item[CodeQL] \textbf{Requires another control}; ---
  \item[Enforcement] Dependabot + dependency review + lockfiles
  \item[Automation] Yes
  \item[Sources] \texttt{NODE-SEC, GH-DEP-REVIEW, GH-DEPENDABOT}
\end{description}
\end{rulebox}
\subsubsection*{Representative implementation example}
\subsubsection*{Example --- assertions are not validation}

\textbf{Noncompliant}\par
\begin{lstlisting}
type User = { id: string };
const user = JSON.parse(body) as User;
useUser(user);
\end{lstlisting}

\textbf{Compliant}\par
\begin{lstlisting}
type User = { id: string };

function isUser(v: unknown): v is User {
  return typeof v === "object" && v !== null &&
         "id" in v && typeof (v as { id?: unknown }).id === "string";
}

const value: unknown = JSON.parse(body);
if (!isUser(value)) throw new Error("Invalid payload");
useUser(value);
\end{lstlisting}

TypeScript documents that type assertions are removed by the compiler and perform no runtime check. Untrusted wire data therefore requires runtime validation.\par
\subsubsection*{Recommended verification controls}
CodeQL JavaScript/TypeScript + \code{tsc} with \code{strict} + runtime schema validation tests + web DAST/IAST as applicable + dependency/secret controls.
\clearpage
\subsection{Python}
Python is memory-safe at the language level but highly dynamic. \code{pickle}, \code{marshal}, \code{eval}/\code{exec}, subprocess/shell use, archive extraction, package management, web frameworks, and native extensions are important trust boundaries. Primary authority: derived from Python 3.14 language/stdlib documentation and official security warnings.
\begin{callout}{Teal}
\textbf{Authority classification:} Derived Secure Coding Baseline\\
\textbf{Primary authority:} Python language/stdlib security warnings and API docs\\
\textbf{Secondary authorities:} NIST/OWASP/MITRE + GitHub
\end{callout}
\subsubsection*{Language-specific secure coding requirements}

\begin{rulebox}[title={\texttt{SCS-PY-DESER-001} \hfill \prioritybadge{Critical Control}{Critical}}]
\textbf{Domain.} Serialization / Deserialization

\textbf{Requirement.} pickle and marshal MUST NOT deserialize untrusted data. Use a data-only format such as JSON plus explicit schema/domain validation when crossing trust boundaries.

\textbf{Rationale.} Python documentation explicitly warns that pickle is not secure and malicious pickle data can execute arbitrary code; marshal is likewise not intended for untrusted data.

\textbf{Risk.} Unsafe deserialization / RCE

\begin{description}[leftmargin=2.65cm,labelwidth=2.3cm,style=nextline,itemsep=1pt,topsep=3pt]
  \item[CWE] CWE-502
  \item[OWASP] OWASP ASVS 5.0 --- deserialization
  \item[CERT] Not mapped
  \item[CodeQL] \textbf{Direct}; py/unsafe-deserialization
  \item[Enforcement] CodeQL + API ban + schema validation
  \item[Automation] Yes
  \item[Sources] \texttt{PY-PICKLE, PY-MARSHAL, GH-BUILTIN}
\end{description}
\end{rulebox}

\begin{rulebox}[title={\texttt{SCS-PY-INJ-001} \hfill \prioritybadge{Critical Control}{Critical}}]
\textbf{Domain.} Injection Prevention

\textbf{Requirement.} subprocess calls with untrusted arguments MUST avoid shell interpretation and pass arguments as a sequence; shell=True MUST NOT be used with untrusted content unless a platform-specific, reviewed escaping design is required.

\textbf{Rationale.} Python does not invoke a shell implicitly, but explicitly enabling a shell transfers quoting/injection responsibility to the application.

\textbf{Risk.} Command injection

\begin{description}[leftmargin=2.65cm,labelwidth=2.3cm,style=nextline,itemsep=1pt,topsep=3pt]
  \item[CWE] CWE-78, CWE-88
  \item[OWASP] Context-dependent / not forced
  \item[CERT] Not mapped
  \item[CodeQL] \textbf{Direct}; py/command-line-injection
  \item[Enforcement] CodeQL + subprocess API policy
  \item[Automation] Yes
  \item[Sources] \texttt{PY-SUBPROC, GH-BUILTIN}
\end{description}
\end{rulebox}

\begin{rulebox}[title={\texttt{SCS-PY-CRYPTO-001} \hfill \prioritybadge{Critical Control}{Critical}}]
\textbf{Domain.} Cryptography

\textbf{Requirement.} Security tokens, passwords-reset values, salts/nonces generated by the application, and similar unpredictable values MUST use the secrets module or an approved cryptographic library, not random.

\textbf{Rationale.} The secrets module is designed for cryptographically strong randomness.

\textbf{Risk.} Predictable security token

\begin{description}[leftmargin=2.65cm,labelwidth=2.3cm,style=nextline,itemsep=1pt,topsep=3pt]
  \item[CWE] CWE-338
  \item[OWASP] Context-dependent / not forced
  \item[CERT] Not mapped
  \item[CodeQL] \textbf{Direct/Partial}; Python insecure-randomness queries
  \item[Enforcement] CodeQL + API policy
  \item[Automation] Partial
  \item[Sources] \texttt{PY-SECRETS, GH-BUILTIN}
\end{description}
\end{rulebox}

\begin{rulebox}[title={\texttt{SCS-PY-FILE-001} \hfill \prioritybadge{High Priority}{High}}]
\textbf{Domain.} File / Path Security

\textbf{Requirement.} Archive extraction of untrusted tar files MUST use the current safe extraction-filter behavior and application-specific inspection/limits; code MUST NOT assume an extraction filter makes arbitrary archives harmless.

\textbf{Rationale.} Archive members can abuse paths, links, device-like metadata, or resource consumption; Python 3.14 changed the default filter to the safer \code{data} behavior but still requires care.

\textbf{Risk.} Path traversal / archive abuse

\begin{description}[leftmargin=2.65cm,labelwidth=2.3cm,style=nextline,itemsep=1pt,topsep=3pt]
  \item[CWE] CWE-22
  \item[OWASP] OWASP ASVS 5.0 --- file handling
  \item[CERT] Not mapped
  \item[CodeQL] \textbf{Direct/Partial}; Python tar/archive extraction and path queries
  \item[Enforcement] CodeQL + extraction policy + tests
  \item[Automation] Partial
  \item[Sources] \texttt{PY-TAR, GH-BUILTIN}
\end{description}
\end{rulebox}

\begin{rulebox}[title={\texttt{SCS-PY-DYN-001} \hfill \prioritybadge{Critical Control}{Critical}}]
\textbf{Domain.} Dynamic Code Execution

\textbf{Requirement.} Untrusted data MUST NOT reach eval, exec, compile-to-code paths, dynamic import names, or template/code evaluation sinks.

\textbf{Rationale.} Dynamic execution directly converts data to executable behavior.

\textbf{Risk.} Code injection

\begin{description}[leftmargin=2.65cm,labelwidth=2.3cm,style=nextline,itemsep=1pt,topsep=3pt]
  \item[CWE] CWE-94, CWE-95
  \item[OWASP] Context-dependent / not forced
  \item[CERT] Not mapped
  \item[CodeQL] \textbf{Direct}; Python code-injection queries
  \item[Enforcement] CodeQL + API restrictions
  \item[Automation] Yes/Partial
  \item[Sources] \texttt{GH-BUILTIN, MITRE-CWE}
\end{description}
\end{rulebox}

\begin{rulebox}[title={\texttt{SCS-PY-BUILD-001} \hfill \prioritybadge{High Priority}{High}}]
\textbf{Domain.} Compiler / Build Security

\textbf{Requirement.} CI MUST pin/reproduce dependency resolution according to the project's packaging policy and SHOULD combine CodeQL with tests, package-vulnerability review, and framework-specific security checks.

\textbf{Rationale.} Python's dynamic behavior and dependency ecosystem require controls beyond static analysis.

\textbf{Risk.} Supply-chain/runtime defects

\begin{description}[leftmargin=2.65cm,labelwidth=2.3cm,style=nextline,itemsep=1pt,topsep=3pt]
  \item[CWE] CWE-829
  \item[OWASP] Context-dependent / not forced
  \item[CERT] Not mapped
  \item[CodeQL] \textbf{Requires another control}; ---
  \item[Enforcement] CodeQL + dependency review/Dependabot + test suite
  \item[Automation] Yes/Partial
  \item[Sources] \texttt{NIST-SSDF, GH-DEP-REVIEW, GH-DEPENDABOT}
\end{description}
\end{rulebox}
\subsubsection*{Representative implementation example}
\subsubsection*{Example --- untrusted serialization}

\textbf{Noncompliant}\par
\begin{lstlisting}[language={Python}]
obj = pickle.loads(request.body)
\end{lstlisting}

\textbf{Compliant}\par
\begin{lstlisting}[language={Python}]
obj = json.loads(request.body)
validate_payload(obj)
\end{lstlisting}

Python explicitly warns that \code{pickle} is not secure and malicious pickle data can execute arbitrary code during unpickling. JSON avoids arbitrary Python object construction, but semantic/schema validation is still required. Relevant weakness: CWE-502. Representative CodeQL query: \code{py/unsafe-deserialization}.\par
\subsubsection*{Recommended verification controls}
CodeQL Python + unit/property/fuzz tests as appropriate + package/dependency review + framework security checks + DAST/IAST for web apps + secret scanning.
\clearpage
\subsection{Ruby}
Ruby is dynamic and includes powerful metaprogramming, process, serialization, and framework facilities. \code{Marshal}, shell-like process invocation, Rails configuration, file/path handling, and gem dependencies require explicit controls. Primary authority: Ruby core/stdlib docs plus Rails Security Guide.
\begin{callout}{Teal}
\textbf{Authority classification:} Derived Secure Coding Baseline\\
\textbf{Primary authority:} Ruby core/stdlib security docs + Rails Security Guide\\
\textbf{Secondary authorities:} NIST/OWASP/MITRE + GitHub
\end{callout}
\subsubsection*{Language-specific secure coding requirements}

\begin{rulebox}[title={\texttt{SCS-RUBY-DESER-001} \hfill \prioritybadge{Critical Control}{Critical}}]
\textbf{Domain.} Serialization / Deserialization

\textbf{Requirement.} Marshal.load MUST NOT process untrusted data; data-only formats and safe loaders such as JSON parsing or appropriately restricted YAML loading MUST be used across trust boundaries.

\textbf{Rationale.} Ruby documentation warns that Marshal can deserialize almost any class and can lead to remote code execution.

\textbf{Risk.} Unsafe deserialization

\begin{description}[leftmargin=2.65cm,labelwidth=2.3cm,style=nextline,itemsep=1pt,topsep=3pt]
  \item[CWE] CWE-502
  \item[OWASP] OWASP ASVS 5.0 --- deserialization
  \item[CERT] Not mapped
  \item[CodeQL] \textbf{Direct}; rb/unsafe-deserialization
  \item[Enforcement] CodeQL + API restrictions
  \item[Automation] Yes
  \item[Sources] \texttt{RUBY-MARSHAL, GH-BUILTIN}
\end{description}
\end{rulebox}

\begin{rulebox}[title={\texttt{SCS-RUBY-INJ-001} \hfill \prioritybadge{Critical Control}{Critical}}]
\textbf{Domain.} Injection Prevention

\textbf{Requirement.} Process execution MUST pass an executable and arguments separately; untrusted data MUST NOT be interpolated into a shell command string.

\textbf{Rationale.} Ruby Open3/Process APIs document command-injection risk with untrusted command strings.

\textbf{Risk.} Command injection

\begin{description}[leftmargin=2.65cm,labelwidth=2.3cm,style=nextline,itemsep=1pt,topsep=3pt]
  \item[CWE] CWE-78, CWE-88
  \item[OWASP] Context-dependent / not forced
  \item[CERT] Not mapped
  \item[CodeQL] \textbf{Direct}; rb/command-line-injection
  \item[Enforcement] CodeQL + Open3/Process API policy
  \item[Automation] Yes
  \item[Sources] \texttt{RUBY-OPEN3, GH-BUILTIN}
\end{description}
\end{rulebox}

\begin{rulebox}[title={\texttt{SCS-RUBY-WEB-001} \hfill \prioritybadge{Critical Control}{Critical}}]
\textbf{Domain.} Web Application Security

\textbf{Requirement.} Rails applications MUST retain and correctly configure framework CSRF, session/cookie, output-escaping, parameter/mass-assignment, and SQL protections rather than bypassing them with raw unsafe APIs.

\textbf{Rationale.} Rails provides security mechanisms, but bypass APIs and unsafe parameter handling can reintroduce common web vulnerabilities.

\textbf{Risk.} XSS/CSRF/SQL/mass assignment

\begin{description}[leftmargin=2.65cm,labelwidth=2.3cm,style=nextline,itemsep=1pt,topsep=3pt]
  \item[CWE] CWE-79, CWE-352, CWE-89, CWE-915
  \item[OWASP] OWASP ASVS 5.0 --- web controls
  \item[CERT] Not mapped
  \item[CodeQL] \textbf{Direct/Partial}; Ruby CodeQL XSS/CSRF/SQL/mass-assignment queries
  \item[Enforcement] CodeQL + Rails config/tests + DAST
  \item[Automation] Partial
  \item[Sources] \texttt{RAILS-SEC, GH-BUILTIN}
\end{description}
\end{rulebox}

\begin{rulebox}[title={\texttt{SCS-RUBY-FILE-001} \hfill \prioritybadge{High Priority}{High}}]
\textbf{Domain.} File / Path Security

\textbf{Requirement.} Untrusted path-like strings MUST be handled by file APIs that cannot reinterpret them as commands/URLs and MUST be constrained to the intended filesystem root.

\textbf{Rationale.} Some Ruby APIs have historically accepted multiple input schemes or shell-like semantics that are unsafe for raw untrusted values.

\textbf{Risk.} Path/command interpretation

\begin{description}[leftmargin=2.65cm,labelwidth=2.3cm,style=nextline,itemsep=1pt,topsep=3pt]
  \item[CWE] CWE-22, CWE-78
  \item[OWASP] Context-dependent / not forced
  \item[CERT] Not mapped
  \item[CodeQL] \textbf{Direct/Partial}; Ruby path and unsafe-open queries
  \item[Enforcement] CodeQL + path policy + tests
  \item[Automation] Partial
  \item[Sources] \texttt{GH-BUILTIN, RUBY-OPEN3}
\end{description}
\end{rulebox}

\begin{rulebox}[title={\texttt{SCS-RUBY-DEP-001} \hfill \prioritybadge{High Priority}{High}}]
\textbf{Domain.} Dependency / Package Security

\textbf{Requirement.} Gem dependencies MUST use reproducible lock state for deployable applications and known-vulnerable gems MUST be reviewed/remediated through GitHub dependency controls or equivalent ecosystem tooling.

\textbf{Rationale.} Dependency compromise can execute within the application/build context.

\textbf{Risk.} Known-vulnerable dependency

\begin{description}[leftmargin=2.65cm,labelwidth=2.3cm,style=nextline,itemsep=1pt,topsep=3pt]
  \item[CWE] CWE-829
  \item[OWASP] Context-dependent / not forced
  \item[CERT] Not mapped
  \item[CodeQL] \textbf{Requires another control}; ---
  \item[Enforcement] Dependabot + dependency review + Bundler lockfile
  \item[Automation] Yes
  \item[Sources] \texttt{GH-DEP-REVIEW, GH-DEPENDABOT, NIST-SSDF}
\end{description}
\end{rulebox}
\subsubsection*{Representative implementation example}
\subsubsection*{Example --- avoid shell command strings}

\textbf{Noncompliant}\par
\begin{lstlisting}[language={Ruby}]
stdout, status = Open3.capture2("git show #{user_ref}")
\end{lstlisting}

\textbf{Compliant}\par
\begin{lstlisting}[language={Ruby}]
stdout, status = Open3.capture2("git", "show", user_ref)
\end{lstlisting}

Ruby's Open3 documentation calls out command-injection risk for untrusted input. Separating the executable and arguments avoids shell expansion in the normal case; application-level validation of allowed refs is still required. Relevant weakness: CWE-78.\par
\subsubsection*{Recommended verification controls}
CodeQL Ruby + Rails security configuration/tests where applicable + dependency review/Dependabot + web DAST/IAST + secret scanning.
\clearpage
\subsection{Rust}
Safe Rust provides a strong soundness property: Safe Rust clients cannot cause undefined behavior. \code{unsafe}, FFI, raw pointers, foreign global state, and external crates form explicit trust boundaries. Memory safety does not prevent SQL injection, SSRF, authorization failures, DoS, or supply-chain compromise. Primary authority: Rustonomicon and Cargo/standard documentation.
\begin{callout}{Teal}
\textbf{Authority classification:} Derived Secure Coding Baseline\\
\textbf{Primary authority:} Rustonomicon + standard/Cargo documentation\\
\textbf{Secondary authorities:} NIST/OWASP/MITRE + GitHub
\end{callout}
\subsubsection*{Language-specific secure coding requirements}

\begin{rulebox}[title={\texttt{SCS-RUST-MEM-001} \hfill \prioritybadge{Critical Control}{Critical}}]
\textbf{Domain.} Memory Safety

\textbf{Requirement.} Application code SHOULD remain in Safe Rust; \code{unsafe} operations MUST be minimized, isolated, documented with required invariants, and reviewed by engineers capable of validating those invariants.

\textbf{Rationale.} Safe Rust clients cannot cause undefined behavior through Safe Rust alone, while unsafe blocks introduce contracts the compiler cannot verify.

\textbf{Risk.} Undefined behavior / memory corruption

\begin{description}[leftmargin=2.65cm,labelwidth=2.3cm,style=nextline,itemsep=1pt,topsep=3pt]
  \item[CWE] CWE-119, CWE-825
  \item[OWASP] Context-dependent / not forced
  \item[CERT] Not mapped
  \item[CodeQL] \textbf{Direct/Partial}; rust/access-after-lifetime-ended (security-extended for representative lifetime case)
  \item[Enforcement] CodeQL + unsafe-code review + tests
  \item[Automation] Partial
  \item[Sources] \texttt{RUST-UNSAFE, GH-BUILTIN}
\end{description}
\end{rulebox}

\begin{rulebox}[title={\texttt{SCS-RUST-FFI-001} \hfill \prioritybadge{Critical Control}{Critical}}]
\textbf{Domain.} Native / Foreign Interfaces

\textbf{Requirement.} Foreign functions MUST be treated as unsafe and SHOULD be wrapped in a small safe Rust interface that validates pointer, length, ownership, error, and thread-safety contracts.

\textbf{Rationale.} Rust cannot verify foreign ABI declarations or the validity/lifetime of foreign pointers.

\textbf{Risk.} Memory corruption / FFI contract violation

\begin{description}[leftmargin=2.65cm,labelwidth=2.3cm,style=nextline,itemsep=1pt,topsep=3pt]
  \item[CWE] CWE-119, CWE-704
  \item[OWASP] Context-dependent / not forced
  \item[CERT] Not mapped
  \item[CodeQL] \textbf{Partial}; ---
  \item[Enforcement] Safe wrapper review + native CodeQL + tests
  \item[Automation] Partial
  \item[Sources] \texttt{RUST-FFI, RUST-UNSAFE}
\end{description}
\end{rulebox}

\begin{rulebox}[title={\texttt{SCS-RUST-RESOURCE-001} \hfill \prioritybadge{High Priority}{High}}]
\textbf{Domain.} Resource Management

\textbf{Requirement.} Untrusted lengths/counts MUST be checked before allocation, collection growth, decompression, parsing, or concurrency fan-out even when memory safety is preserved.

\textbf{Rationale.} Memory-safe code can still exhaust memory/CPU or trigger integer/resource logic defects.

\textbf{Risk.} Resource exhaustion

\begin{description}[leftmargin=2.65cm,labelwidth=2.3cm,style=nextline,itemsep=1pt,topsep=3pt]
  \item[CWE] CWE-400, CWE-789
  \item[OWASP] Context-dependent / not forced
  \item[CERT] Not mapped
  \item[CodeQL] \textbf{Direct/Partial}; Rust allocation-size and regex/resource queries
  \item[Enforcement] CodeQL + fuzz/load tests + limits
  \item[Automation] Partial
  \item[Sources] \texttt{GH-BUILTIN, NIST-SSDF}
\end{description}
\end{rulebox}

\begin{rulebox}[title={\texttt{SCS-RUST-NET-001} \hfill \prioritybadge{Critical Control}{Critical}}]
\textbf{Domain.} Network Communications

\textbf{Requirement.} Outbound requests built from untrusted URLs MUST enforce approved schemes/hosts/networks and prevent access to protected internal destinations where SSRF is relevant.

\textbf{Rationale.} Memory safety does not prevent server-side request forgery.

\textbf{Risk.} SSRF

\begin{description}[leftmargin=2.65cm,labelwidth=2.3cm,style=nextline,itemsep=1pt,topsep=3pt]
  \item[CWE] CWE-918
  \item[OWASP] OWASP ASVS 5.0 --- SSRF
  \item[CERT] Not mapped
  \item[CodeQL] \textbf{Direct}; rust/request-forgery
  \item[Enforcement] CodeQL + egress policy + tests
  \item[Automation] Yes/Partial
  \item[Sources] \texttt{GH-BUILTIN, OWASP-ASVS}
\end{description}
\end{rulebox}

\begin{rulebox}[title={\texttt{SCS-RUST-DEP-001} \hfill \prioritybadge{High Priority}{High}}]
\textbf{Domain.} Dependency / Package Security

\textbf{Requirement.} Deployable applications MUST commit and honor Cargo.lock according to Cargo guidance and MUST review vulnerable crates/transitive dependencies.

\textbf{Rationale.} Locking gives deterministic dependency selection; vulnerability review addresses known security defects.

\textbf{Risk.} Known-vulnerable dependency

\begin{description}[leftmargin=2.65cm,labelwidth=2.3cm,style=nextline,itemsep=1pt,topsep=3pt]
  \item[CWE] CWE-829
  \item[OWASP] Context-dependent / not forced
  \item[CERT] Not mapped
  \item[CodeQL] \textbf{Requires another control}; ---
  \item[Enforcement] Cargo.lock + Dependabot/dependency review + ecosystem audit tooling
  \item[Automation] Yes
  \item[Sources] \texttt{CARGO-LOCK, GH-DEP-REVIEW, GH-DEPENDABOT}
\end{description}
\end{rulebox}

\begin{rulebox}[title={\texttt{SCS-RUST-INJ-001} \hfill \prioritybadge{Critical Control}{Critical}}]
\textbf{Domain.} Injection Prevention

\textbf{Requirement.} Safe Rust's memory guarantees MUST NOT be treated as injection prevention; SQL, command, template, and URL sinks MUST still use structured APIs and validation.

\textbf{Rationale.} Memory safety does not constrain interpreter semantics.

\textbf{Risk.} SQL/command/template injection

\begin{description}[leftmargin=2.65cm,labelwidth=2.3cm,style=nextline,itemsep=1pt,topsep=3pt]
  \item[CWE] CWE-89, CWE-78, CWE-94
  \item[OWASP] OWASP ASVS 5.0 --- injection
  \item[CERT] Not mapped
  \item[CodeQL] \textbf{Direct/Partial}; rust/sql-injection plus other Rust injection queries
  \item[Enforcement] CodeQL + structured APIs
  \item[Automation] Partial
  \item[Sources] \texttt{GH-BUILTIN, OWASP-ASVS}
\end{description}
\end{rulebox}
\subsubsection*{Representative implementation example}
\subsubsection*{Example --- isolate FFI unsafety}

\textbf{Preferred pattern}\par
\begin{lstlisting}
unsafe extern "C" {
    fn foreign_check(ptr: *const u8, len: usize) -> i32;
}

pub fn check(data: &[u8]) -> bool {
    // SAFETY: data.as_ptr() is valid for data.len() bytes for this call;
    // the foreign function does not retain the pointer.
    unsafe { foreign_check(data.as_ptr(), data.len()) == 0 }
}
\end{lstlisting}

The unsafe operation is kept inside a small safe wrapper whose contract is documented. The FFI declaration and foreign implementation still require independent validation. Rust's documentation states that FFI calls are unsafe because the compiler cannot verify the foreign side.\par
\subsubsection*{Recommended verification controls}
CodeQL Rust + compiler/lints + focused review of every \code{unsafe}/FFI boundary + fuzz/property tests + dependency review + secret scanning.
\clearpage
\subsection{Swift}
Safe Swift provides initialization, bounds, lifetime, and overlapping-access/thread-safety protections, while intentionally exposing unsafe pointer and C/C++ interop features. On Apple platforms, CryptoKit, Keychain, ATS/URLSession, WebKit, and platform entitlements/configuration are key security surfaces. Primary authority: Swift language docs plus Apple platform security docs.
\begin{callout}{Teal}
\textbf{Authority classification:} Derived Secure Coding Baseline\\
\textbf{Primary authority:} Swift language safety docs + Apple platform security docs\\
\textbf{Secondary authorities:} NIST/OWASP/MITRE + GitHub
\end{callout}
\subsubsection*{Language-specific secure coding requirements}

\begin{rulebox}[title={\texttt{SCS-SWIFT-MEM-001} \hfill \prioritybadge{Critical Control}{Critical}}]
\textbf{Domain.} Memory Safety

\textbf{Requirement.} Code SHOULD use safe Swift collections/references and MUST isolate UnsafePointer/UnsafeMutablePointer or other unchecked memory operations behind small reviewed interfaces.

\textbf{Rationale.} Swift provides initialization, bounds, lifetime, and overlap/thread-safety protections for safe code but also exposes intentionally unsafe pointer constructs.

\textbf{Risk.} Memory corruption / lifetime error

\begin{description}[leftmargin=2.65cm,labelwidth=2.3cm,style=nextline,itemsep=1pt,topsep=3pt]
  \item[CWE] CWE-119, CWE-416
  \item[OWASP] Context-dependent / not forced
  \item[CERT] Not mapped
  \item[CodeQL] \textbf{Partial}; Swift unsafe-memory/pointer-related queries where available
  \item[Enforcement] CodeQL + compiler/runtime checks + review
  \item[Automation] Partial
  \item[Sources] \texttt{SWIFT-MEM, SWIFT-INTEROP, GH-BUILTIN}
\end{description}
\end{rulebox}

\begin{rulebox}[title={\texttt{SCS-SWIFT-FFI-001} \hfill \prioritybadge{Critical Control}{Critical}}]
\textbf{Domain.} Native / Foreign Interfaces

\textbf{Requirement.} C/C++ interoperability MUST document and validate lifetime, ownership, nullability, pointer bounds, and thread-safety assumptions; unsafe references MUST NOT outlive the underlying foreign object/storage.

\textbf{Rationale.} Interoperability can escape Swift's normal memory-safety model.

\textbf{Risk.} Memory/lifetime violation

\begin{description}[leftmargin=2.65cm,labelwidth=2.3cm,style=nextline,itemsep=1pt,topsep=3pt]
  \item[CWE] CWE-119, CWE-825
  \item[OWASP] Context-dependent / not forced
  \item[CERT] Not mapped
  \item[CodeQL] \textbf{Partial}; ---
  \item[Enforcement] Interop review + Swift/native CodeQL + tests
  \item[Automation] Partial
  \item[Sources] \texttt{SWIFT-INTEROP}
\end{description}
\end{rulebox}

\begin{rulebox}[title={\texttt{SCS-SWIFT-CRYPTO-001} \hfill \prioritybadge{Critical Control}{Critical}}]
\textbf{Domain.} Cryptography

\textbf{Requirement.} Apple-platform applications SHOULD use CryptoKit high-level APIs for supported cryptographic operations and MUST NOT implement custom primitives or raw-pointer cryptography without specialized review.

\textbf{Rationale.} Apple recommends CryptoKit over lower-level interfaces; it handles security-sensitive memory details for common operations.

\textbf{Risk.} Weak/custom cryptography

\begin{description}[leftmargin=2.65cm,labelwidth=2.3cm,style=nextline,itemsep=1pt,topsep=3pt]
  \item[CWE] CWE-327
  \item[OWASP] OWASP ASVS 5.0 --- cryptography
  \item[CERT] Not mapped
  \item[CodeQL] \textbf{Direct/Partial}; Swift weak-crypto/static-IV/iteration-count queries
  \item[Enforcement] CodeQL + cryptographic review
  \item[Automation] Partial
  \item[Sources] \texttt{APPLE-CRYPTO, GH-BUILTIN}
\end{description}
\end{rulebox}

\begin{rulebox}[title={\texttt{SCS-SWIFT-SECRET-001} \hfill \prioritybadge{Critical Control}{Critical}}]
\textbf{Domain.} Secrets

\textbf{Requirement.} Small application secrets such as credentials, private keys, and tokens SHOULD be stored using Keychain services with appropriate access-control attributes rather than plaintext files/preferences.

\textbf{Rationale.} Keychain provides encrypted storage and platform access-control mechanisms for secrets.

\textbf{Risk.} Sensitive data exposure

\begin{description}[leftmargin=2.65cm,labelwidth=2.3cm,style=nextline,itemsep=1pt,topsep=3pt]
  \item[CWE] CWE-312
  \item[OWASP] OWASP ASVS 5.0 --- secrets
  \item[CERT] Not mapped
  \item[CodeQL] \textbf{Partial}; Swift cleartext-storage/logging queries
  \item[Enforcement] CodeQL + Keychain + secret scanning
  \item[Automation] Partial
  \item[Sources] \texttt{APPLE-KEYCHAIN, GH-SECRET, GH-BUILTIN}
\end{description}
\end{rulebox}

\begin{rulebox}[title={\texttt{SCS-SWIFT-NET-001} \hfill \prioritybadge{Critical Control}{Critical}}]
\textbf{Domain.} Network Communications

\textbf{Requirement.} Apple-platform applications SHOULD rely on URLSession/ATS defaults for HTTPS and MUST NOT globally disable ATS or trust evaluation for production traffic without a documented exception and compensating controls.

\textbf{Rationale.} ATS requires TLS and reliable certificates/ciphers for URL Loading System traffic; broad exceptions weaken this protection.

\textbf{Risk.} Insecure TLS / cleartext transport

\begin{description}[leftmargin=2.65cm,labelwidth=2.3cm,style=nextline,itemsep=1pt,topsep=3pt]
  \item[CWE] CWE-295, CWE-319, CWE-757
  \item[OWASP] OWASP ASVS 5.0 --- communications
  \item[CERT] Not mapped
  \item[CodeQL] \textbf{Direct}; swift/insecure-tls plus cleartext-transport queries
  \item[Enforcement] CodeQL + ATS configuration review + integration tests
  \item[Automation] Yes/Partial
  \item[Sources] \texttt{APPLE-ATS, GH-BUILTIN}
\end{description}
\end{rulebox}

\begin{rulebox}[title={\texttt{SCS-SWIFT-WEB-001} \hfill \prioritybadge{High Priority}{High}}]
\textbf{Domain.} Web Application Security

\textbf{Requirement.} WKWebView/JavaScriptCore integrations MUST treat loaded content, navigation targets, message handlers, and JavaScript bridging as trust boundaries and restrict them to approved origins/capabilities.

\textbf{Rationale.} Embedding web/script runtimes introduces DOM/navigation/code-execution risks not eliminated by Swift type/memory safety.

\textbf{Risk.} Web content / script injection

\begin{description}[leftmargin=2.65cm,labelwidth=2.3cm,style=nextline,itemsep=1pt,topsep=3pt]
  \item[CWE] CWE-79, CWE-94
  \item[OWASP] Context-dependent / not forced
  \item[CERT] Not mapped
  \item[CodeQL] \textbf{Direct/Partial}; Swift unsafe WebView query family
  \item[Enforcement] CodeQL + navigation/content policy + tests
  \item[Automation] Partial
  \item[Sources] \texttt{GH-BUILTIN, OWASP-ASVS}
\end{description}
\end{rulebox}
\subsubsection*{Representative implementation example}
\subsubsection*{Example --- preserve ATS rather than disabling it}

\textbf{Preferred}\par
\begin{lstlisting}
let url = URL(string: "https://api.example.com/v1/items")!
let (data, response) = try await URLSession.shared.data(from: url)
\end{lstlisting}

Do not add a global \code{NSAllowsArbitraryLoads = YES} merely to make an HTTP/TLS problem disappear. Apple documents that ATS is enforced by the URL Loading System and that broad exceptions weaken its protections. Network authorization, response validation, and application-level trust checks remain required.\par
\subsubsection*{Recommended verification controls}
CodeQL Swift on macOS + compiler/concurrency diagnostics + platform security/configuration review + unit/UI/network tests + dependency/secret controls + interop review.


---
\clearpage
\section{GitHub Actions Secure Workflow Standard}
GitHub Actions is treated as a separate CodeQL analysis target rather than as a programming language. Its security model is a CI/CD control-plane model: attacker-controlled repository metadata, pull-request code, workflow expressions, tokens, secrets, third-party actions, caches/artifacts, runners, and cloud credentials cross trust and privilege boundaries.

\begin{rulebox}[title={\texttt{SCS-ACTIONS-INJ-001} \hfill \prioritybadge{Critical Control}{Critical}}]
\textbf{Domain.} Injection Prevention

\textbf{Requirement.} Untrusted GitHub context values MUST NOT be interpolated directly into \code{run:} scripts; pass them through an intermediate environment variable and quote/use them as data for the target shell.

\textbf{Rationale.} Expression expansion happens before the shell executes and can turn pull-request metadata or other attacker-controlled fields into shell syntax.

\textbf{Risk.} Code / command injection

\begin{description}[leftmargin=2.65cm,labelwidth=2.3cm,style=nextline,itemsep=1pt,topsep=3pt]
  \item[CWE] CWE-94, CWE-95, CWE-116
  \item[OWASP] Context-dependent / not forced
  \item[CERT] Not mapped
  \item[CodeQL] \textbf{Direct}; actions/code-injection/critical
  \item[Enforcement] CodeQL required check + workflow review
  \item[Automation] Yes
  \item[Sources] \texttt{GH-SECURE-ACTIONS, GH-BUILTIN}
\end{description}
\end{rulebox}

\begin{rulebox}[title={\texttt{SCS-ACTIONS-AUTHZ-001} \hfill \prioritybadge{Critical Control}{Critical}}]
\textbf{Domain.} Authorization

\textbf{Requirement.} Each workflow MUST declare the minimum \code{GITHUB\_TOKEN} permissions required; write permissions MUST be granted only to jobs that need them.

\textbf{Rationale.} Over-privileged workflow tokens increase the blast radius of workflow or dependency compromise.

\textbf{Risk.} Excess privilege

\begin{description}[leftmargin=2.65cm,labelwidth=2.3cm,style=nextline,itemsep=1pt,topsep=3pt]
  \item[CWE] CWE-275, CWE-250
  \item[OWASP] Context-dependent / not forced
  \item[CERT] Not mapped
  \item[CodeQL] \textbf{Direct}; actions/missing-workflow-permissions
  \item[Enforcement] CodeQL + organization workflow policy
  \item[Automation] Yes
  \item[Sources] \texttt{GH-SECURE-ACTIONS, GH-BUILTIN}
\end{description}
\end{rulebox}

\begin{rulebox}[title={\texttt{SCS-ACTIONS-DEP-001} \hfill \prioritybadge{Critical Control}{Critical}}]
\textbf{Domain.} Dependency / Package Security

\textbf{Requirement.} Third-party actions SHOULD be pinned to a full-length commit SHA; tag-only references require explicit risk acceptance and compensating review.

\textbf{Rationale.} Tags can move; a full commit SHA is immutable and reduces dependency substitution risk.

\textbf{Risk.} Untrusted dependency / action

\begin{description}[leftmargin=2.65cm,labelwidth=2.3cm,style=nextline,itemsep=1pt,topsep=3pt]
  \item[CWE] CWE-829
  \item[OWASP] Context-dependent / not forced
  \item[CERT] Not mapped
  \item[CodeQL] \textbf{Direct}; actions/unpinned-tag (security-extended/security-and-quality)
  \item[Enforcement] CodeQL security-extended + Actions policy + Dependabot
  \item[Automation] Yes
  \item[Sources] \texttt{GH-SECURE-ACTIONS, GH-BUILTIN}
\end{description}
\end{rulebox}

\begin{rulebox}[title={\texttt{SCS-ACTIONS-SECRET-001} \hfill \prioritybadge{Critical Control}{Critical}}]
\textbf{Domain.} Secrets

\textbf{Requirement.} Workflow secrets MUST be minimized, scoped to the smallest job/environment, never echoed, and MUST NOT be exposed to untrusted fork/pull-request code; cloud authentication SHOULD use OIDC short-lived credentials where supported.

\textbf{Rationale.} Workflow logs, subprocesses, third-party actions, and untrusted code can exfiltrate long-lived secrets.

\textbf{Risk.} Secret exposure / credential theft

\begin{description}[leftmargin=2.65cm,labelwidth=2.3cm,style=nextline,itemsep=1pt,topsep=3pt]
  \item[CWE] CWE-312, CWE-532
  \item[OWASP] Context-dependent / not forced
  \item[CERT] Not mapped
  \item[CodeQL] \textbf{Direct/Partial}; Actions excessive-secret-exposure and unmasked-secret queries
  \item[Enforcement] CodeQL + secret scanning + environments + OIDC
  \item[Automation] Partial
  \item[Sources] \texttt{GH-SECURE-ACTIONS, GH-OIDC, GH-SECRET, GH-BUILTIN}
\end{description}
\end{rulebox}

\begin{rulebox}[title={\texttt{SCS-ACTIONS-AUTHZ-002} \hfill \prioritybadge{Critical Control}{Critical}}]
\textbf{Domain.} Authorization

\textbf{Requirement.} \code{pull\_request\_target} or other privileged triggers MUST NOT check out and execute untrusted pull-request code with write tokens or secrets.

\textbf{Rationale.} Privileged base-repository context combined with attacker-controlled code creates a repository compromise path.

\textbf{Risk.} Untrusted code with elevated privileges

\begin{description}[leftmargin=2.65cm,labelwidth=2.3cm,style=nextline,itemsep=1pt,topsep=3pt]
  \item[CWE] CWE-829, CWE-269
  \item[OWASP] Context-dependent / not forced
  \item[CERT] Not mapped
  \item[CodeQL] \textbf{Direct/Partial}; Actions untrusted-checkout/privileged-trigger queries
  \item[Enforcement] CodeQL + workflow architecture review
  \item[Automation] Partial
  \item[Sources] \texttt{GH-SECURE-ACTIONS, GH-BUILTIN}
\end{description}
\end{rulebox}

\begin{rulebox}[title={\texttt{SCS-ACTIONS-BUILD-001} \hfill \prioritybadge{High Priority}{High}}]
\textbf{Domain.} Compiler / Build Security

\textbf{Requirement.} Security workflows SHOULD use CodeQL default/code-scanning queries plus \code{security-extended} where the organization accepts the additional signal, and required checks SHOULD block merges according to documented severity policy.

\textbf{Rationale.} Some security queries, including unpinned-action detection and representative Rust lifetime analysis, are not in the default code-scanning suite.

\textbf{Risk.} Coverage gap from suite selection

\begin{description}[leftmargin=2.65cm,labelwidth=2.3cm,style=nextline,itemsep=1pt,topsep=3pt]
  \item[CWE] Not forced / context-dependent
  \item[OWASP] Context-dependent / not forced
  \item[CERT] Not mapped
  \item[CodeQL] \textbf{Requires configuration}; Query-suite configuration
  \item[Enforcement] CodeQL workflow + branch/ruleset required checks
  \item[Automation] Yes
  \item[Sources] \texttt{GH-BUILTIN, GH-WORKFLOW}
\end{description}
\end{rulebox}

\begin{rulebox}[title={\texttt{SCS-ACTIONS-FILE-001} \hfill \prioritybadge{High Priority}{High}}]
\textbf{Domain.} File / Path Security

\textbf{Requirement.} Artifacts and caches from untrusted or lower-trust jobs MUST NOT be consumed by privileged jobs as executable/configuration input without integrity and provenance controls.

\textbf{Rationale.} Artifact/cache poisoning can move attacker-controlled content across job or privilege boundaries.

\textbf{Risk.} Artifact/cache poisoning

\begin{description}[leftmargin=2.65cm,labelwidth=2.3cm,style=nextline,itemsep=1pt,topsep=3pt]
  \item[CWE] CWE-829, CWE-94
  \item[OWASP] Context-dependent / not forced
  \item[CERT] Not mapped
  \item[CodeQL] \textbf{Direct/Partial}; Actions artifact-poisoning / cache-poisoning query families
  \item[Enforcement] CodeQL + workflow isolation + provenance checks
  \item[Automation] Partial
  \item[Sources] \texttt{GH-SECURE-ACTIONS, GH-BUILTIN}
\end{description}
\end{rulebox}

\begin{rulebox}[title={\texttt{SCS-ACTIONS-DEP-002} \hfill \prioritybadge{High Priority}{High}}]
\textbf{Domain.} Dependency / Package Security

\textbf{Requirement.} Pull requests that change dependency manifests/lockfiles or action references SHOULD run dependency review and MUST block additions that violate the organization's vulnerability/license policy.

\textbf{Rationale.} Dependency review exposes dependency changes at the pull-request boundary, before merge.

\textbf{Risk.} Vulnerable dependency introduction

\begin{description}[leftmargin=2.65cm,labelwidth=2.3cm,style=nextline,itemsep=1pt,topsep=3pt]
  \item[CWE] CWE-829
  \item[OWASP] Context-dependent / not forced
  \item[CERT] Not mapped
  \item[CodeQL] \textbf{Requires another control}; Dependency Review Action
  \item[Enforcement] Dependency review required check + Dependabot
  \item[Automation] Yes
  \item[Sources] \texttt{GH-DEP-REVIEW, GH-DEPENDABOT}
\end{description}
\end{rulebox}
\subsection*{Representative secure workflow pattern}
\subsubsection*{Example --- untrusted context in shell}

\textbf{Noncompliant}\par
\begin{lstlisting}
- name: Print PR title
  run: echo "${{ github.event.pull_request.title }}"
\end{lstlisting}

\textbf{Compliant}\par
\begin{lstlisting}
- name: Print PR title
  env:
    PR_TITLE: ${{ github.event.pull_request.title }}
  run: printf '%s\n' "$PR_TITLE"
\end{lstlisting}

GitHub recommends using an intermediate environment variable for untrusted context values so expression expansion does not splice attacker-controlled syntax directly into the generated script. Representative CodeQL query: \code{actions/code-injection/critical}.\par

---\par
\clearpage
\begin{landscape}
\section{Verified Representative CodeQL Query Crosswalk}
This table records query metadata preserved in the source catalog. It is representative rather than exhaustive; GitHub maintains a larger built-in query set for each supported language.

\scriptsize
\setlength{\tabcolsep}{3pt}
\begin{longtable}{@{}P{0.9in}P{1.25in}P{1.85in}P{1.0in}P{0.7in}P{0.55in}P{2.15in}@{}}
\toprule
\textbf{Language} & \textbf{Weakness / query} & \textbf{Query ID} & \textbf{CWE} & \textbf{Security severity} & \textbf{Precision} & \textbf{Suites}\\
\midrule
\endfirsthead
\toprule
\textbf{Language} & \textbf{Weakness / query} & \textbf{Query ID} & \textbf{CWE} & \textbf{Security severity} & \textbf{Precision} & \textbf{Suites}\\
\midrule
\endhead
C/C++ & SQL injection & cpp/sql-injection & CWE-89 & 8.8 & high & code-scanning; security-extended; security-and-quality\\
C/C++ & Command-line injection & cpp/command-line-injection & CWE-78, CWE-88 & 9.8 & high & code-scanning; security-extended; security-and-quality\\
C/C++ & Use after free & cpp/use-after-free & CWE-416 & 9.3 & high & code-scanning; security-extended; security-and-quality\\
C/C++ & Weak cryptographic algorithm & cpp/weak-cryptographic-algorithm & CWE-327 & 7.5 & high & code-scanning; security-extended; security-and-quality\\
C\# & SQL injection & cs/sql-injection & CWE-89 & 8.8 & high & code-scanning; security-extended; security-and-quality\\
C\# & Unsafe deserialization of untrusted data & cs/unsafe-deserialization-untrusted-input & CWE-502 & 9.8 & high & code-scanning; security-extended; security-and-quality\\
C\# & Command-line injection & cs/command-line-injection & CWE-78, CWE-88 & 9.8 & high & code-scanning; security-extended; security-and-quality\\
Go & SQL injection & go/sql-injection & CWE-89 & 8.8 & high & code-scanning; security-extended; security-and-quality\\
Go & Server-side request forgery & go/request-forgery & CWE-918 & 9.1 & high & code-scanning; security-extended; security-and-quality\\
Go & Disabled certificate check & go/disabled-certificate-check & CWE-295 & 7.5 & high & code-scanning; security-extended; security-and-quality\\
Java/Kotlin & SQL injection & java/sql-injection & CWE-89, CWE-564 & 8.8 & high & code-scanning; security-extended; security-and-quality\\
Java/Kotlin & Unsafe deserialization & java/unsafe-deserialization & CWE-502 & 9.8 & high & code-scanning; security-extended; security-and-quality\\
Java/Kotlin & Command-line injection & java/command-line-injection & CWE-78, CWE-88 & 9.8 & high & code-scanning; security-extended; security-and-quality\\
JavaScript/TypeScript & SQL injection & js/sql-injection & CWE-89, CWE-90, CWE-943 & 8.8 & high & code-scanning; security-extended; security-and-quality\\
JavaScript/TypeScript & Prototype-polluting assignment & js/prototype-polluting-assignment & CWE-78, 79, 94, 400, 471, 915 & 6.1 & high & code-scanning; security-extended; security-and-quality\\
JavaScript/TypeScript & Server-side request forgery & js/request-forgery & CWE-918 & 9.1 & high & code-scanning; security-extended; security-and-quality\\
Python & SQL injection & py/sql-injection & CWE-89 & 8.8 & high & code-scanning; security-extended; security-and-quality\\
Python & Unsafe deserialization & py/unsafe-deserialization & CWE-502 & 9.8 & high & code-scanning; security-extended; security-and-quality\\
Python & Command-line injection & py/command-line-injection & CWE-78, CWE-88 & 9.8 & high & code-scanning; security-extended; security-and-quality\\
Ruby & SQL injection & rb/sql-injection & CWE-89 & 8.8 & high & code-scanning; security-extended; security-and-quality\\
Ruby & Unsafe deserialization & rb/unsafe-deserialization & CWE-502 & 9.8 & high & code-scanning; security-extended; security-and-quality\\
Ruby & Command-line injection & rb/command-line-injection & CWE-78, CWE-88 & 9.8 & high & code-scanning; security-extended; security-and-quality\\
Rust & SQL injection & rust/sql-injection & CWE-89 & 8.8 & high & code-scanning; security-extended; security-and-quality\\
Rust & Access after lifetime ended & rust/access-after-lifetime-ended & CWE-825 & 9.8 & medium & security-extended; security-and-quality\\
Rust & Server-side request forgery & rust/request-forgery & CWE-918 & 9.1 & high & code-scanning; security-extended; security-and-quality\\
Swift & SQL injection & swift/sql-injection & CWE-89 & 8.8 & high & code-scanning; security-extended; security-and-quality\\
Swift & Command-line injection & swift/command-line-injection & CWE-78, CWE-88 & 9.8 & high & code-scanning; security-extended; security-and-quality\\
Swift & Insecure TLS & swift/insecure-tls & CWE-757 & 7.5 & high & code-scanning; security-extended; security-and-quality\\
GitHub Actions & Code injection & actions/code-injection/critical & CWE-94, CWE-95, CWE-116 & 9.0 & very-high & code-scanning; security-extended; security-and-quality\\
GitHub Actions & Missing workflow permissions & actions/missing-workflow-permissions & CWE-275 & 5.0 & high & code-scanning; security-extended; security-and-quality\\
GitHub Actions & Unpinned action tag & actions/unpinned-tag & CWE-829 & 5.0 & medium & security-extended; security-and-quality\\
\bottomrule
\end{longtable}
\end{landscape}

\subsection{How to Read Detection Coverage}
\begin{tabularx}{\textwidth}{@{}P{1.7in}Y@{}}
\toprule
\textbf{Coverage class} & \textbf{Interpretation}\\
\midrule
Direct & A verified CodeQL query or query family can directly identify representative instances of the weakness pattern.\\
Partial & CodeQL can identify important instances, but coverage depends on APIs, frameworks, data flow, build/extraction, or the exact coding pattern.\\
Not reliably detectable & Static analysis is not a dependable primary verification method for the requirement.\\
Requires another control & Use another GitHub security capability, language-native tool, dynamic test, runtime control, or governance mechanism as the primary enforcement method.\\
\bottomrule
\end{tabularx}
\clearpage
\begin{landscape}
\section{Master Traceability Matrix}
The traceability matrix connects each normative rule to its weakness taxonomy, CodeQL coverage, representative query/control, and enforcement method.

\tiny
\setlength{\tabcolsep}{2.5pt}
\begin{longtable}{@{}P{1.0in}P{0.6in}P{0.95in}P{0.75in}P{0.85in}P{0.75in}P{2.0in}P{2.05in}@{}}
\toprule
\textbf{Rule ID} & \textbf{Lang.} & \textbf{Domain} & \textbf{Priority} & \textbf{CWE} & \textbf{CodeQL} & \textbf{Query / control} & \textbf{Enforcement}\\
\midrule
\endfirsthead
\toprule
\textbf{Rule ID} & \textbf{Lang.} & \textbf{Domain} & \textbf{Priority} & \textbf{CWE} & \textbf{CodeQL} & \textbf{Query / control} & \textbf{Enforcement}\\
\midrule
\endhead
SCS-ALL-INPUT-001 & All & Input Validation & High Priority & CWE-20 & Partial & Language-specific taint/data-flow queries & CodeQL + schema/parser validation + unit/fuzz tests\\
SCS-ALL-OUTPUT-001 & All & Output Encoding & High Priority & CWE-79, CWE-116 & Partial & Language/framework-specific XSS and response-splitting queries & CodeQL + framework-safe templates + tests\\
SCS-ALL-INJ-001 & All & Injection Prevention & Critical Control & CWE-78, CWE-89, CWE-94, CWE-95, CWE-943 & Direct & Representative: cpp/sql-injection; cs/sql-injection; go/sql-injection; java/sql-injection; js/sql-injection; py/sql-injection; rb/sql-injection; rust/sql-injection; swift/sql-injection & CodeQL required check + API-level parameterization\\
SCS-ALL-MEM-001 & All & Memory Safety & Critical Control & CWE-119, CWE-416, CWE-415, CWE-825 & Partial & Representative: cpp/use-after-free; rust/access-after-lifetime-ended & CodeQL + compiler/runtime checks + sanitizers/dynamic tests + review of unsafe/native boundaries\\
SCS-ALL-AUTHN-001 & All & Authentication & Critical Control & CWE-287, CWE-306 & Partial & Framework-specific authentication queries & Architecture review + framework configuration tests + DAST/IAST where applicable\\
SCS-ALL-AUTHZ-001 & All & Authorization & Critical Control & CWE-285, CWE-862, CWE-863, CWE-639 & Partial & Language/framework-specific missing authorization / IDOR queries & CodeQL where modeled + authorization tests + architecture/code review\\
SCS-ALL-CRYPTO-001 & All & Cryptography & Critical Control & CWE-327, CWE-326, CWE-338 & Partial & Weak-crypto / insecure-randomness / TLS queries vary by language & CodeQL + platform API policy + cryptographic design review\\
SCS-ALL-SECRET-001 & All & Secrets & Critical Control & CWE-798, CWE-532, CWE-312 & Requires another control & GitHub secret scanning / push protection rather than CodeQL for primary enforcement & Secret scanning + push protection + approved secret store + rotation process\\
SCS-ALL-ERROR-001 & All & Error / Exception Handling & High Priority & CWE-209, CWE-703 & Partial & Information-disclosure and unchecked-error queries vary & Code review + negative-path tests + logging policy + CodeQL where available\\
SCS-ALL-LOG-001 & All & Logging / Monitoring & High Priority & CWE-117, CWE-532 & Partial & Log-injection / cleartext-logging queries exist for multiple languages & CodeQL + centralized logging policy + tests\\
SCS-ALL-FILE-001 & All & File / Path Security & High Priority & CWE-22, CWE-59, CWE-367, CWE-377 & Direct/Partial & Path-injection/traversal, archive extraction, temporary-file, and TOCTOU queries vary by language & CodeQL + filesystem permission tests + archive policy\\
SCS-ALL-DESER-001 & All & Serialization / Deserialization & Critical Control & CWE-502 & Direct/Partial & Representative: cs/unsafe-deserialization-untrusted-input; java/unsafe-deserialization; py/unsafe-deserialization; rb/unsafe-deserialization & CodeQL + API allowlist + schema validation\\
SCS-ALL-NET-001 & All & Network Communications & Critical Control & CWE-295, CWE-297, CWE-319, CWE-757 & Direct/Partial & Representative: go/disabled-certificate-check; swift/insecure-tls; TLS/certificate queries in other packs & CodeQL + platform network policy + integration tests\\
SCS-ALL-WEB-001 & All & Web Application Security & Critical Control & CWE-79, CWE-352, CWE-918, CWE-601 & Direct/Partial & Representative: go/request-forgery; js/request-forgery; rust/request-forgery; framework-specific XSS/CSRF queries & CodeQL + framework config tests + DAST/IAST + architecture review\\
SCS-ALL-DEP-001 & All & Dependency / Package Security & Critical Control & CWE-829 & Requires another control & Dependabot alerts/security updates + dependency review & Dependabot + dependency review + lockfiles + approved registries\\
SCS-ALL-RESOURCE-001 & All & Resource Management & High Priority & CWE-400, CWE-770, CWE-789, CWE-1333 & Partial & Allocation-size, regex, decompression and other resource queries vary by language & CodeQL + fuzz/load tests + runtime quotas/timeouts\\
SCS-ALL-CONC-001 & All & Concurrency / Parallelism & High Priority & CWE-362, CWE-367 & Partial & TOCTOU/static queries where available; dynamic race detectors for Go and native code & CodeQL + race detector/thread sanitizer where supported + concurrency tests\\
SCS-ALL-DYN-001 & All & Dynamic Code Execution & Critical Control & CWE-94, CWE-95 & Direct/Partial & Code-injection / dynamic-code queries vary by language & CodeQL + code review + API restrictions\\
SCS-ALL-FFI-001 & All & Native / Foreign Interfaces & Critical Control & CWE-119, CWE-416, CWE-704 & Partial & Native-side CodeQL where source is analyzed; limited cross-language guarantees & Boundary review + native/static analysis + sanitizers + tests\\
SCS-ALL-BUILD-001 & All & Compiler / Build Security & High Priority & CWE-1104 (where applicable) & Requires another control & CodeQL check runs plus language-native compiler/linter/test/dependency checks & CI/CD policy + CodeQL + ecosystem-native tooling\\
SCS-C-MEM-001 & C & Memory Safety & Critical Control & CWE-119, CWE-787 & Direct/Partial & C/C++ bounds and buffer-overflow queries & CodeQL + compiler diagnostics + sanitizers + review\\
SCS-C-MEM-002 & C & Memory Safety & Critical Control & CWE-416, CWE-415 & Direct & cpp/use-after-free and related double-free queries & CodeQL + ASan-capable test builds + review\\
SCS-C-INPUT-001 & C & Input Validation & High Priority & CWE-190, CWE-681 & Partial & C/C++ overflow and conversion queries & CodeQL + compiler diagnostics + boundary tests\\
SCS-C-OUTPUT-001 & C & Output Encoding & High Priority & CWE-134 & Direct & C/C++ uncontrolled-format-string query & CodeQL + compiler format warnings\\
SCS-C-FILE-001 & C & File / Path Security & High Priority & CWE-367 & Direct/Partial & C/C++ TOCTOU queries & CodeQL + handle-based APIs + tests\\
SCS-C-BUILD-001 & C & Compiler / Build Security & High Priority & --- & Requires another control & --- & Compiler diagnostics + sanitizer test jobs + CodeQL\\
SCS-CPP-MEM-001 & C++ & Memory Safety & Critical Control & CWE-416, CWE-415 & Partial & C/C++ lifetime, mismatched allocation/deallocation, and UAF queries & CodeQL + review + sanitizer tests\\
SCS-CPP-MEM-002 & C++ & Memory Safety & Critical Control & CWE-416, CWE-825 & Direct/Partial & cpp/use-after-free and iterator/lifetime queries & CodeQL + compiler/sanitizer + review\\
SCS-CPP-INPUT-001 & C++ & Input Validation & High Priority & CWE-190, CWE-681 & Partial & C/C++ integer and allocation-size queries & CodeQL + compiler conversion warnings + tests\\
SCS-CPP-FFI-001 & C++ & Native / Foreign Interfaces & Critical Control & CWE-119, CWE-704 & Partial & --- & Wrapper review + CodeQL + sanitizer tests\\
SCS-CPP-CONC-001 & C++ & Concurrency / Parallelism & High Priority & CWE-362 & Partial & --- & Thread sanitizer / tests + review; CodeQL where modeled\\
SCS-CPP-BUILD-001 & C++ & Compiler / Build Security & High Priority & --- & Requires another control & --- & Compiler + sanitizer tests + CodeQL\\
SCS-CSHARP-DESER-001 & C\# & Serialization / Deserialization & Critical Control & CWE-502 & Direct & cs/unsafe-deserialization-untrusted-input & CodeQL + API ban/analyzer + migration to data-only serializers\\
SCS-CSHARP-CRYPTO-001 & C\# & Cryptography & Critical Control & CWE-338 & Direct/Partial & C\# insecure-randomness / weak-crypto queries & CodeQL + API restrictions + tests\\
SCS-CSHARP-FFI-001 & C\# & Native / Foreign Interfaces & Critical Control & CWE-119, CWE-704 & Partial & --- & Code review + native-side analysis + tests\\
SCS-CSHARP-WEB-001 & C\# & Web Application Security & Critical Control & CWE-352 & Partial & ASP.NET CSRF-related CodeQL queries where modeled & Framework configuration + integration/DAST tests + CodeQL\\
SCS-CSHARP-AUTHZ-001 & C\# & Authorization & Critical Control & CWE-862, CWE-639 & Partial & C\# access-control / IDOR queries where framework modeling applies & Authorization tests + CodeQL + review\\
SCS-CSHARP-DYN-001 & C\# & Dynamic Code Execution & Critical Control & CWE-78, CWE-94 & Direct/Partial & cs/command-line-injection plus dynamic-code queries & CodeQL + API restrictions + review\\
SCS-GO-OUTPUT-001 & Go & Output Encoding & High Priority & CWE-79 & Direct/Partial & Go XSS queries & CodeQL + template API policy + tests\\
SCS-GO-INJ-001 & Go & Injection Prevention & Critical Control & CWE-78 & Direct & Go command-injection query & CodeQL + os/exec API policy\\
SCS-GO-NET-001 & Go & Network Communications & Critical Control & CWE-295 & Direct & go/disabled-certificate-check & CodeQL required check + integration tests\\
SCS-GO-CRYPTO-001 & Go & Cryptography & Critical Control & CWE-338 & Direct/Partial & Go insecure-randomness queries & CodeQL + API policy\\
SCS-GO-CONC-001 & Go & Concurrency / Parallelism & High Priority & CWE-362 & Requires another control & --- & go test -race + tests + CodeQL for other issues\\
SCS-GO-DEP-001 & Go & Dependency / Package Security & High Priority & CWE-1104 & Requires another control & --- & govulncheck + Dependabot + dependency review\\
SCS-JAVA-DESER-001 & Java & Serialization / Deserialization & Critical Control & CWE-502 & Direct & java/unsafe-deserialization & CodeQL + serializer allowlist/schema + review\\
SCS-JAVA-INJ-001 & Java & Injection Prevention & Critical Control & CWE-89, CWE-78 & Direct & java/sql-injection; java/command-line-injection & CodeQL + API policy\\
SCS-JAVA-INPUT-001 & Java & Input Validation & High Priority & CWE-20 & Partial & Multiple Java taint/data-flow queries & CodeQL + parser/schema tests\\
SCS-JAVA-RESOURCE-001 & Java & Resource Management & High Priority & CWE-400, CWE-770, CWE-1333 & Partial & Java resource/regex/XML queries & CodeQL + limits + fuzz/load tests\\
SCS-JAVA-FFI-001 & Java & Native / Foreign Interfaces & Critical Control & CWE-119, CWE-704 & Partial & --- & Java + native CodeQL where both sources are available; boundary review\\
SCS-JAVA-DEP-001 & Java & Dependency / Package Security & High Priority & CWE-1104 & Requires another control & --- & Dependency review + Dependabot + ecosystem SBOM/scanner\\
SCS-KOTLIN-INPUT-001 & Kotlin & Input Validation & High Priority & CWE-20 & Partial & Java/Kotlin data-flow queries where modeled & Schema/domain validation + CodeQL\\
SCS-KOTLIN-MEM-001 & Kotlin & Memory Safety & High Priority & CWE-476 & Partial & Java/Kotlin nullness-related queries where applicable & Compiler + review + tests\\
SCS-KOTLIN-FFI-001 & Kotlin & Native / Foreign Interfaces & Critical Control & CWE-704, CWE-119 & Partial & --- & Java/Kotlin CodeQL + boundary review\\
SCS-KOTLIN-DEP-001 & Kotlin & Dependency / Package Security & High Priority & CWE-1104 & Requires another control & --- & Dependabot/dependency review + ecosystem vulnerability scanning\\
SCS-KOTLIN-DESER-001 & Kotlin & Serialization / Deserialization & Critical Control & CWE-502 & Direct/Partial & java/unsafe-deserialization where the modeled JVM API is used & CodeQL + serializer policy\\
SCS-JS-DYN-001 & JavaScript & Dynamic Code Execution & Critical Control & CWE-94, CWE-95 & Direct & JavaScript code-injection queries & CodeQL + lint/API restrictions\\
SCS-JS-INPUT-001 & JavaScript & Input Validation & High Priority & CWE-1321, CWE-915 & Direct & js/prototype-polluting-assignment & CodeQL + safe merge/object patterns + tests\\
SCS-JS-INJ-001 & JavaScript & Injection Prevention & Critical Control & CWE-78 & Direct & JavaScript command-injection queries & CodeQL + child\_process API policy\\
SCS-JS-WEB-001 & JavaScript & Web Application Security & Critical Control & CWE-79 & Direct & JavaScript XSS / unsafe DOM queries & CodeQL + framework-safe rendering + DAST\\
SCS-JS-RESOURCE-001 & JavaScript & Resource Management & High Priority & CWE-400, CWE-1333 & Partial & JavaScript regex/resource queries & CodeQL + runtime limits + load/fuzz tests\\
SCS-JS-DEP-001 & JavaScript & Dependency / Package Security & Critical Control & CWE-829 & Requires another control & --- & Dependabot + dependency review + lockfile-controlled CI\\
SCS-TS-INPUT-001 & TypeScript & Input Validation & High Priority & CWE-20 & Partial & JavaScript/TypeScript data-flow queries & Runtime schema/domain validation + compiler + CodeQL\\
SCS-TS-INPUT-002 & TypeScript & Input Validation & High Priority & CWE-20, CWE-704 & Partial & --- & Compiler/review + runtime validation tests\\
SCS-TS-BUILD-001 & TypeScript & Compiler / Build Security & High Priority & --- & Requires another control & --- & tsc --noEmit/CI + CodeQL\\
SCS-TS-DYN-001 & TypeScript & Dynamic Code Execution & Critical Control & CWE-79, CWE-78, CWE-94 & Direct/Partial & JavaScript/TypeScript injection and prototype-pollution queries & CodeQL + JavaScript runtime controls\\
SCS-TS-DEP-001 & TypeScript & Dependency / Package Security & Critical Control & CWE-829 & Requires another control & --- & Dependabot + dependency review + lockfiles\\
SCS-PY-DESER-001 & Python & Serialization / Deserialization & Critical Control & CWE-502 & Direct & py/unsafe-deserialization & CodeQL + API ban + schema validation\\
SCS-PY-INJ-001 & Python & Injection Prevention & Critical Control & CWE-78, CWE-88 & Direct & py/command-line-injection & CodeQL + subprocess API policy\\
SCS-PY-CRYPTO-001 & Python & Cryptography & Critical Control & CWE-338 & Direct/Partial & Python insecure-randomness queries & CodeQL + API policy\\
SCS-PY-FILE-001 & Python & File / Path Security & High Priority & CWE-22 & Direct/Partial & Python tar/archive extraction and path queries & CodeQL + extraction policy + tests\\
SCS-PY-DYN-001 & Python & Dynamic Code Execution & Critical Control & CWE-94, CWE-95 & Direct & Python code-injection queries & CodeQL + API restrictions\\
SCS-PY-BUILD-001 & Python & Compiler / Build Security & High Priority & CWE-829 & Requires another control & --- & CodeQL + dependency review/Dependabot + test suite\\
SCS-RUBY-DESER-001 & Ruby & Serialization / Deserialization & Critical Control & CWE-502 & Direct & rb/unsafe-deserialization & CodeQL + API restrictions\\
SCS-RUBY-INJ-001 & Ruby & Injection Prevention & Critical Control & CWE-78, CWE-88 & Direct & rb/command-line-injection & CodeQL + Open3/Process API policy\\
SCS-RUBY-WEB-001 & Ruby & Web Application Security & Critical Control & CWE-79, CWE-352, CWE-89, CWE-915 & Direct/Partial & Ruby CodeQL XSS/CSRF/SQL/mass-assignment queries & CodeQL + Rails config/tests + DAST\\
SCS-RUBY-FILE-001 & Ruby & File / Path Security & High Priority & CWE-22, CWE-78 & Direct/Partial & Ruby path and unsafe-open queries & CodeQL + path policy + tests\\
SCS-RUBY-DEP-001 & Ruby & Dependency / Package Security & High Priority & CWE-829 & Requires another control & --- & Dependabot + dependency review + Bundler lockfile\\
SCS-RUST-MEM-001 & Rust & Memory Safety & Critical Control & CWE-119, CWE-825 & Direct/Partial & rust/access-after-lifetime-ended (security-extended for representative lifetime case) & CodeQL + unsafe-code review + tests\\
SCS-RUST-FFI-001 & Rust & Native / Foreign Interfaces & Critical Control & CWE-119, CWE-704 & Partial & --- & Safe wrapper review + native CodeQL + tests\\
SCS-RUST-RESOURCE-001 & Rust & Resource Management & High Priority & CWE-400, CWE-789 & Direct/Partial & Rust allocation-size and regex/resource queries & CodeQL + fuzz/load tests + limits\\
SCS-RUST-NET-001 & Rust & Network Communications & Critical Control & CWE-918 & Direct & rust/request-forgery & CodeQL + egress policy + tests\\
SCS-RUST-DEP-001 & Rust & Dependency / Package Security & High Priority & CWE-829 & Requires another control & --- & Cargo.lock + Dependabot/dependency review + ecosystem audit tooling\\
SCS-RUST-INJ-001 & Rust & Injection Prevention & Critical Control & CWE-89, CWE-78, CWE-94 & Direct/Partial & rust/sql-injection plus other Rust injection queries & CodeQL + structured APIs\\
SCS-SWIFT-MEM-001 & Swift & Memory Safety & Critical Control & CWE-119, CWE-416 & Partial & Swift unsafe-memory/pointer-related queries where available & CodeQL + compiler/runtime checks + review\\
SCS-SWIFT-FFI-001 & Swift & Native / Foreign Interfaces & Critical Control & CWE-119, CWE-825 & Partial & --- & Interop review + Swift/native CodeQL + tests\\
SCS-SWIFT-CRYPTO-001 & Swift & Cryptography & Critical Control & CWE-327 & Direct/Partial & Swift weak-crypto/static-IV/iteration-count queries & CodeQL + cryptographic review\\
SCS-SWIFT-SECRET-001 & Swift & Secrets & Critical Control & CWE-312 & Partial & Swift cleartext-storage/logging queries & CodeQL + Keychain + secret scanning\\
SCS-SWIFT-NET-001 & Swift & Network Communications & Critical Control & CWE-295, CWE-319, CWE-757 & Direct & swift/insecure-tls plus cleartext-transport queries & CodeQL + ATS configuration review + integration tests\\
SCS-SWIFT-WEB-001 & Swift & Web Application Security & High Priority & CWE-79, CWE-94 & Direct/Partial & Swift unsafe WebView query family & CodeQL + navigation/content policy + tests\\
SCS-ACTIONS-INJ-001 & GitHub Actions & Injection Prevention & Critical Control & CWE-94, CWE-95, CWE-116 & Direct & actions/code-injection/critical & CodeQL required check + workflow review\\
SCS-ACTIONS-AUTHZ-001 & GitHub Actions & Authorization & Critical Control & CWE-275, CWE-250 & Direct & actions/missing-workflow-permissions & CodeQL + organization workflow policy\\
SCS-ACTIONS-DEP-001 & GitHub Actions & Dependency / Package Security & Critical Control & CWE-829 & Direct & actions/unpinned-tag (security-extended/security-and-quality) & CodeQL security-extended + Actions policy + Dependabot\\
SCS-ACTIONS-SECRET-001 & GitHub Actions & Secrets & Critical Control & CWE-312, CWE-532 & Direct/Partial & Actions excessive-secret-exposure and unmasked-secret queries & CodeQL + secret scanning + environments + OIDC\\
SCS-ACTIONS-AUTHZ-002 & GitHub Actions & Authorization & Critical Control & CWE-829, CWE-269 & Direct/Partial & Actions untrusted-checkout/privileged-trigger queries & CodeQL + workflow architecture review\\
SCS-ACTIONS-BUILD-001 & GitHub Actions & Compiler / Build Security & High Priority & --- & Requires configuration & Query-suite configuration & CodeQL workflow + branch/ruleset required checks\\
SCS-ACTIONS-FILE-001 & GitHub Actions & File / Path Security & High Priority & CWE-829, CWE-94 & Direct/Partial & Actions artifact-poisoning / cache-poisoning query families & CodeQL + workflow isolation + provenance checks\\
SCS-ACTIONS-DEP-002 & GitHub Actions & Dependency / Package Security & High Priority & CWE-829 & Requires another control & Dependency Review Action & Dependency review required check + Dependabot\\
\bottomrule
\end{longtable}
\end{landscape}
\begin{landscape}
\section{Enterprise Enforcement Matrix}
\scriptsize
\setlength{\tabcolsep}{3pt}
\begin{longtable}{@{}P{1.2in}P{1.75in}P{2.0in}P{1.85in}P{1.65in}@{}}
\toprule
\textbf{Requirement class} & \textbf{Primary prevention} & \textbf{Static / security automation} & \textbf{Dynamic / operational verification} & \textbf{Human / governance control}\\
\midrule
Input / output / injection & Structured APIs, validation, contextual encoding & CodeQL data-flow and security queries & Fuzzing, integration tests, DAST/IAST & Trust-boundary and code review\\
Native memory & Safe abstractions, explicit ownership and bounds & CodeQL C/C++/Rust/Swift queries; compiler diagnostics & ASan/UBSan/TSan where supported; fuzzing & Unsafe/native boundary review\\
AuthN/AuthZ & Framework/IdP policies; default deny & Partial framework-aware CodeQL & Authorization matrix tests; DAST/IAST & Architecture/security review\\
Cryptography & Platform crypto APIs; approved key management & Weak-crypto/TLS/randomness queries & Interoperability/network tests & Cryptographic design review\\
Secrets & External secret store; scoped credentials & Secret scanning; push protection & Rotation/revocation validation & Secret ownership/governance\\
Dependencies & Approved registries; lockfiles/pinning & Dependabot; dependency review; ecosystem scanners & Build/test upgraded dependency & Exception and upgrade-SLA governance\\
Serialization & Data-only formats; restricted loaders & Unsafe-deserialization queries & Malicious payload tests/fuzzing & Serializer/API review\\
Filesystem & Root confinement; handle-based APIs; safe extraction & Path/Zip Slip/temp/TOCTOU queries & Filesystem integration tests & Privilege/filesystem design review\\
Network & TLS defaults; peer verification; outbound allowlists & TLS/certificate/SSRF queries & Integration tests; egress monitoring & Exception review\\
Resource/concurrency & Limits, backpressure, ownership/synchronization & Partial allocation/regex/race queries & Load/fuzz/race-detector testing & Capacity/threat modeling\\
CI/CD / Actions & Least privilege; immutable refs; OIDC; isolation & Actions CodeQL + dependency review & Ephemeral runner/artifact validation & Workflow/ruleset governance\\
\bottomrule
\end{longtable}
\end{landscape}
\normalsize
\clearpage
\section{Adoption Guidance}
\subsection{Minimum Enterprise Baseline}
\begin{enumerate}[leftmargin=*,itemsep=3pt]
  \item Apply every \code{SCS-ALL-*} rule to every repository unless a documented applicability exception exists.
  \item Add the \code{SCS-<LANG>-*} overlays for every language detected in the repository.
  \item Enable CodeQL for every supported language present. Test \code{security-extended} in observation mode before making it merge-blocking, because broader suites can change signal/noise and some valuable queries are not in the default suite.
  \item Enable secret scanning and push protection; treat an exposed credential as an incident requiring rotation or revocation, not merely deletion from the current file.
  \item Enable dependency graph, Dependabot, and dependency review where product configuration and policy allow; make dependency review a required check where the organization's risk policy demands it.
  \item Require least-privilege Actions permissions, immutable third-party action references, restricted privileged triggers, OIDC for supported cloud access, and explicit runner/artifact/cache trust boundaries.
  \item Pair CodeQL with DAST/IAST for web/API behavior, language-native race/sanitizer/fuzz testing, and authorization/business-logic tests.
  \item Record owner, rationale, compensating controls, expiry, and remediation plan for every secure-coding exception.
\end{enumerate}

\subsection{What CodeQL Cannot Prove}
CodeQL does not prove that the business authorization model is correct; that an authentication protocol or configuration is fit for purpose; that cryptographic key lifecycle and operational custody are correct; that secrets are rotated after exposure; that runtime and infrastructure configuration are secure; that all dependencies are trustworthy merely because no known vulnerability is present; or that an application is secure because analysis has no alerts.

\subsection{Fact-Check and Maintenance Notes}
\begin{itemize}[leftmargin=*,itemsep=2pt]
  \item The language/version table is point-in-time and should be reverified when the CodeQL bundle changes.
  \item Framework/library coverage indicates published modeling, not complete semantic coverage of every API.
  \item Exact CWE mapping is intentionally omitted where an architectural rule can lead to several weakness classes.
  \item OWASP ASVS mappings are retained at the control-area level unless an exact requirement has been separately verified.
  \item Some CodeQL queries live only in \code{security-extended} or \code{security-and-quality}; suite selection is therefore a material security-policy decision.
  \item Managed or memory-safe languages do not remove logic, injection, authorization, cryptographic, dependency, or resource-exhaustion risks.
\end{itemize}
\clearpage
\begin{landscape}
\section{Official Source Catalog}
All source entries are first-party or authoritative references preserved from the source catalog. Source keys used in rule cards resolve here.

\scriptsize
\setlength{\tabcolsep}{3pt}
\begin{longtable}{@{}P{0.82in}P{1.23in}P{2.25in}P{1.7in}P{3.0in}@{}}
\toprule
\textbf{Key} & \textbf{Organization} & \textbf{Document} & \textbf{Applies to} & \textbf{URL}\\
\midrule
\endfirsthead
\toprule
\textbf{Key} & \textbf{Organization} & \textbf{Document} & \textbf{Applies to} & \textbf{URL}\\
\midrule
\endhead
APPLE-ATS & Apple & Preventing Insecure Network Connections & Swift / Apple network transport & \url{https://developer.apple.com/documentation/security/preventing-insecure-network-connections}\\
APPLE-CRYPTO & Apple & Apple CryptoKit & Swift / Apple cryptography & \url{https://developer.apple.com/documentation/cryptokit}\\
APPLE-KEYCHAIN & Apple & Keychain services & Swift / Apple secret storage & \url{https://developer.apple.com/documentation/security/keychain-services}\\
ASPNET-CSRF & Microsoft & Prevent Cross-Site Request Forgery (XSRF/CSRF) attacks in ASP.NET Core & C\# / ASP.NET Core & \url{https://learn.microsoft.com/en-us/aspnet/core/security/anti-request-forgery}\\
CARGO-LOCK & The Rust Project & Cargo Guide --- Cargo.toml vs Cargo.lock & Rust dependency locking & \url{https://doc.rust-lang.org/cargo/guide/cargo-toml-vs-cargo-lock.html}\\
CERT-C & SEI / Carnegie Mellon University & SEI CERT C Coding Standard, 2016 Edition & C & \url{https://www.sei.cmu.edu/library/sei-cert-c-coding-standard-rules-for-developing-safe-reliable-and-secure-systems-2016-edition/}\\
CERT-CPP & SEI / Carnegie Mellon University & SEI CERT C++ Coding Standard, 2016 Edition & C++ & \url{https://www.sei.cmu.edu/library/sei-cert-c-coding-standard-rules-for-developing-safe-reliable-and-secure-systems-2016-edition-2/}\\
CERT-JAVA & SEI / Carnegie Mellon University & SEI CERT Oracle Coding Standard for Java & Java secure-coding rules; published standard scope is older Java SE, with development pages continuing to receive updates & \url{https://wiki.sei.cmu.edu/confluence/display/java}\\
DOTNET-BF & Microsoft & Deserialization risks in use of BinaryFormatter and related types & C\# / .NET deserialization & \url{https://learn.microsoft.com/dotnet/standard/serialization/binaryformatter-security-guide}\\
DOTNET-CRYPTO & Microsoft & Cryptographic services & C\# / .NET cryptography & \url{https://learn.microsoft.com/en-us/dotnet/standard/security/cryptographic-services}\\
DOTNET-SEC & Microsoft & Secure coding guidelines for .NET & C\# / .NET & \url{https://learn.microsoft.com/en-us/dotnet/standard/security/secure-coding-guidelines}\\
GH-BUILTIN & GitHub & CodeQL built-in queries & Built-in CodeQL query inventories & \url{https://docs.github.com/en/code-security/reference/code-scanning/codeql/codeql-queries/built-in-queries}\\
GH-DEP-REVIEW & GitHub & About dependency review & Dependency review and pull-request gating & \url{https://docs.github.com/en/code-security/supply-chain-security/understanding-your-software-supply-chain/about-dependency-review}\\
GH-DEPENDABOT & GitHub & About Dependabot security updates & Vulnerable dependency remediation & \url{https://docs.github.com/en/code-security/dependabot/dependabot-security-updates/about-dependabot-security-updates}\\
GH-LANG & GitHub & CodeQL --- Supported languages and frameworks & All CodeQL languages; versions, compilers, frameworks, libraries & \url{https://codeql.github.com/docs/codeql-overview/supported-languages-and-frameworks/}\\
GH-OIDC & GitHub & About security hardening with OpenID Connect & Short-lived cloud credentials from Actions & \url{https://docs.github.com/en/actions/concepts/security/openid-connect}\\
GH-PUSH & GitHub & About push protection & Blocking supported secrets before repository exposure & \url{https://docs.github.com/en/code-security/secret-scanning/introduction/about-push-protection}\\
GH-SECRET & GitHub & About secret scanning & Secret discovery & \url{https://docs.github.com/en/code-security/secret-scanning/introduction/about-secret-scanning}\\
GH-SECURE-ACTIONS & GitHub & Secure use reference --- GitHub Actions & GitHub Actions workflow security & \url{https://docs.github.com/en/actions/reference/security/secure-use}\\
GH-WORKFLOW & GitHub & Configuring CodeQL workflow / language identifiers and build modes & CodeQL language identifiers and compiled-language build behavior & \url{https://docs.github.com/en/code-security/code-scanning/creating-an-advanced-setup-for-code-scanning/configuring-codeql-workflow-for-compiled-languages}\\
GO-EXEC & The Go Project & os/exec package & Go process execution & \url{https://pkg.go.dev/os/exec}\\
GO-RAND & The Go Project & crypto/rand package & Go cryptographic randomness & \url{https://pkg.go.dev/crypto/rand}\\
GO-SEC & The Go Project & Security Best Practices for Go Developers & Go & \url{https://go.dev/doc/security/best-practices}\\
GO-TEMPLATE & The Go Project & html/template package & Go output encoding / HTML templating & \url{https://pkg.go.dev/html/template}\\
JAVA-26 & Oracle & JDK 26 Documentation --- Security & Java 26 security APIs and specifications & \url{https://docs.oracle.com/en/java/javase/26/}\\
JAVA-SCG & Oracle & Secure Coding Guidelines for Java SE, document version 11.0 (updated June 2025) & Java & \url{https://www.oracle.com/java/technologies/javase/seccodeguide.html}\\
KOTLIN-NULL & JetBrains & Kotlin Null safety & Kotlin nullability and Java interoperability & \url{https://kotlinlang.org/docs/null-safety.html}\\
KOTLIN-SEC & JetBrains & Kotlin Security & Kotlin & \url{https://kotlinlang.org/docs/security.html}\\
MITRE-CWE & MITRE & Common Weakness Enumeration (CWE) & Weakness taxonomy and mappings & \url{https://cwe.mitre.org/}\\
NIST-SSDF & NIST & SP 800-218, Secure Software Development Framework (SSDF) Version 1.1 & Cross-language secure software development practices & \url{https://csrc.nist.gov/pubs/sp/800/218/final}\\
NODE-CHILD & OpenJS Foundation / Node.js & Node.js child\_process API & JavaScript process execution & \url{https://nodejs.org/api/child_process.html}\\
NODE-SEC & OpenJS Foundation / Node.js & Node.js Security Best Practices & Server-side JavaScript / Node.js & \url{https://nodejs.org/en/learn/getting-started/security-best-practices}\\
OWASP-ASVS & OWASP & Application Security Verification Standard (ASVS) 5.0 & Web/application security verification requirements & \url{https://owasp.org/www-project-application-security-verification-standard/}\\
PY-MARSHAL & Python Software Foundation & Python 3.14 --- marshal & Python deserialization & \url{https://docs.python.org/3.14/library/marshal.html}\\
PY-PICKLE & Python Software Foundation & Python 3.14 --- pickle & Python deserialization & \url{https://docs.python.org/3.14/library/pickle.html}\\
PY-SECRETS & Python Software Foundation & Python 3.14 --- secrets & Python cryptographic randomness & \url{https://docs.python.org/3.14/library/secrets.html}\\
PY-SUBPROC & Python Software Foundation & Python 3.14 --- subprocess security considerations & Python command execution & \url{https://docs.python.org/3.14/library/subprocess.html\#security-considerations}\\
PY-TAR & Python Software Foundation & Python 3.14 --- tarfile & Python archive extraction & \url{https://docs.python.org/3.14/library/tarfile.html}\\
RAILS-SEC & Ruby on Rails & Ruby on Rails Security Guide & Ruby / Rails web security & \url{https://guides.rubyonrails.org/security.html}\\
RUBY-MARSHAL & Ruby Core & Ruby 3.3 Marshal --- Security considerations & Ruby deserialization & \url{https://ruby-doc.org/3.3/core/Marshal.html}\\
RUBY-OPEN3 & Ruby Core & Ruby 3.3 Open3 & Ruby process execution & \url{https://ruby-doc.org/3.3/stdlibs/open3/Open3.html}\\
RUST-FFI & The Rust Project & The Rustonomicon --- FFI & Rust foreign-function interfaces & \url{https://doc.rust-lang.org/nomicon/ffi.html}\\
RUST-UNSAFE & The Rust Project & The Rustonomicon --- How Safe and Unsafe Interact & Rust unsafe boundaries & \url{https://doc.rust-lang.org/stable/nomicon/safe-unsafe-meaning.html}\\
SWIFT-INTEROP & The Swift Project & Safely Mixing Swift and C/C++ & Swift C/C++ interoperability & \url{https://www.swift.org/documentation/cxx-interop/safe-interop/}\\
SWIFT-MEM & The Swift Project & The Swift Programming Language --- The Basics: Memory Safety & Swift memory safety & \url{https://docs.swift.org/swift-book/documentation/the-swift-programming-language/thebasics/}\\
TS-STRICT & Microsoft / TypeScript & TSConfig --- strict & TypeScript strict type checking & \url{https://www.typescriptlang.org/tsconfig/strict}\\
TS-TYPES & Microsoft / TypeScript & TypeScript Handbook --- Everyday Types & TypeScript any, assertions, runtime limitations & \url{https://www.typescriptlang.org/docs/handbook/2/everyday-types.html}\\
\bottomrule
\end{longtable}
\end{landscape}
\clearpage
\appendix
\section{Rule Identifier Scheme}
\begin{callout}{Teal}
\centering
\Large\code{SCS-<LANG>-<DOMAIN>-<NNN>}
\end{callout}

\textbf{Language identifiers} include \code{C}, \code{CPP}, \code{CSHARP}, \code{GO}, \code{JAVA}, \code{KOTLIN}, \code{JS}, \code{TS}, \code{PY}, \code{RUBY}, \code{RUST}, \code{SWIFT}, \code{ACTIONS}, and \code{ALL}.

\textbf{Domain identifiers} include \code{INPUT}, \code{OUTPUT}, \code{INJ}, \code{MEM}, \code{AUTHN}, \code{AUTHZ}, \code{CRYPTO}, \code{SECRET}, \code{ERROR}, \code{LOG}, \code{FILE}, \code{DESER}, \code{NET}, \code{WEB}, \code{DEP}, \code{RESOURCE}, \code{CONC}, \code{DYN}, \code{FFI}, and \code{BUILD}.

\section{Normative Interpretation}
\begin{description}[leftmargin=2.0cm,labelwidth=1.65cm,style=nextline]
  \item[MUST / MUST NOT] Mandatory enterprise baseline unless an approved exception exists.
  \item[SHOULD / SHOULD NOT] Expected practice; deviation requires documented engineering rationale.
  \item[MAY] Optional technique or control depending on context.
\end{description}

This catalog intentionally separates \textbf{prevention}, \textbf{detection}, and \textbf{verification}. A CodeQL-detectable pattern can still require runtime prevention, and a non-CodeQL-detectable rule remains a security requirement.

\end{document}
